\chapter{Extracting the Pauli observables}
\label{chap:qpbt-extraction}
% Paper subsection label sec:separating kept for cross-reference fidelity.
\label{sec:separating}
% Source: references/qpbt-paper/14_analysis_of_the_pauli_basis_test.tex,
% lines 1415--1877 (subsection ``Pulling the X and Z measurements apart''
% and the closing proof of the soundness theorem).

This chapter completes the analysis of the Pauli basis test~\cite{Ji2020MIPStar}. Throughout, $(q,m,d)$ is an admissible parameter tuple (Definition~\ref{def:admissible}), $q = 2^t$, and $M = 2^m$ denotes the number of $q$-ary qudits tested for per player. (The letter $M$ also denotes measurement operators, always carrying super- and subscripts; the meaning is determined by context, as in the source.) We fix a projective strategy for the Pauli basis test game $\mathfrak{G}^{\mathrm{Pauli}}_{(q,m,d)}$ that succeeds with probability at least $1-\eps$, in the setting of Chapter~\ref{chap:qpbt-observables}: the shared state $\ket{\psi}$ on registers $\mathrm{A}$, $\mathrm{B}$ is assumed padded with sufficiently many ancilla qubits in the state $\ket{0}$ (Lemma~\ref{lem:projective-strategy-setup} and the convention following it), and $\ket{\hat\psi}$ denotes the expanded state on the six registers $\mathrm{A}\mathrm{A}'\mathrm{A}''\mathrm{B}\mathrm{B}'\mathrm{B}''$, obtained by adjoining to each player the two halves of $M$ pairs $\ket{\mathrm{EPR}_q}$ (Definition~\ref{def:expanded-state}). Unless stated otherwise, every $\approx_\delta$ and $\simeq_\delta$ statement below is taken with respect to the state $\ket{\hat\psi}$, over the question distribution named in the statement, with register placements indicated by subscripts; ``all symmetric equivalents'' refers to the register interchanges of Definition~\ref{def:symmetric-equivalents}.

We use the following objects from the preceding chapters: the combined point projectors $\hat{M}^{(\mathrm{Point},W),u}_a$ acting on the expanded space (Definition~\ref{def:expanded-point-measurement}); the projective polynomial-pair measurement $\{\hat{S}_{g_X,g_Z}\}$ of Lemma~\ref{lem:qld-4-7}, with outcomes pairs of polynomials in $\mathrm{ideg}_{d,m}(\F_q)$ and error function
\[
	\delta_G(\eps,m,d,q) = a(md)^a\bigl(\eps^b + q^{-b} + 2^{-bmd}\bigr)
\]
for universal constants $a > 1$, $0 < b < 1$. This is the error denoted $\delta_S$ in Lemma~\ref{lem:qld-4-7}; we call it $\delta_G$ here to distinguish the global polynomial-pair witness error from the enlarged error used below. We also use the decoding map $\mathrm{Dec} = \mathrm{Dec}_{\F_q}\colon \mathrm{ideg}_{d,m}(\F_q) \to \F_q^{M}$ of the low-degree code (Definition~\ref{def:decoding-map}, applied with $H = \F_q$; see Remark~\ref{rem:qld-decoding-identity}), which maps a polynomial $g$ to the string of its evaluations on the hypercube $\{0,1\}^m$; and the low-degree encoding $h \mapsto g_h$ together with the vectors $\mathrm{ind}_m(u) \in \F_q^M$, characterized by $h \cdot \mathrm{ind}_m(u) = g_h(u)$ for all $u \in \F_q^m$ \eqref{eq:low-degree-encoding-definition}. Both maps are $\F_q$-linear, and $\mathrm{Dec}(g_h) = h$ for every $h \in \F_q^{M}$. We call a polynomial $g \in \mathrm{ideg}_{d,m}(\F_q)$ an \emph{encoding} if $g = g_h$ for some $h \in \F_q^{M}$ or, equivalently, if $g$ is multilinear; since a multilinear polynomial is determined by its values on $\{0,1\}^m$, the encodings are exactly the $g$ with $g_{\mathrm{Dec}(g)} = g$ as polynomial representatives, and
\[
	\mathrm{Dec}(g) \cdot \mathrm{ind}_m(u) = g(u)
	\qquad \text{for every encoding $g$ and every } u \in \F_q^m\;.
\]
For a $g$ that is not an encoding the two sides of this identity are the evaluations of the distinct polynomial representatives $g_{\mathrm{Dec}(g)}$ and $g$; when $d \leq q-1$ they consequently disagree at some $u$, but when $d \geq q$ the two representatives may induce the same function on $\F_q^m$, so pointwise validity of the identity does not characterize the encodings; see Remark~\ref{rem:qld-decoding-identity}.

\begin{definition}[Full-field decoding on bounded polynomials]\label{def:qld-full-field-decoder}
	\lean{MIPStarRE.QPBT.evalPoly, MIPStarRE.QPBT.decodeOnUniv, MIPStarRE.QPBT.decodeFq}
	\lean{MIPStarRE.QPBT.encodingPoly, MIPStarRE.QPBT.IsEncoding}
	\leanok
	\uses{def:decoding-map, def:low-degree-encoding}
	For a bounded polynomial representative $g$, let $\mathrm{Dec}_{\F_q}(g)$ be its evaluations on the Boolean cube, with every value of $\F_q$ retained. For $h \in \F_q^M$, let $g_h$ denote its multilinear encoding. The predicate $\mathrm{IsEncoding}(g)$ asserts the equality $g_{\mathrm{Dec}_{\F_q}(g)}=g$ of polynomial representatives.
\end{definition}

\begin{lemma}[Linearity and inversion of the full-field decoder]\label{lem:qld-decoder-linearity}
	\lean{MIPStarRE.QPBT.decodeFq_add, MIPStarRE.QPBT.decodeFq_smul, MIPStarRE.QPBT.decodeFq_lowDegreeEncoding}
	\leanok
	\uses{def:qld-full-field-decoder}
	For bounded polynomial representatives $g,g'$ and $c\in\F_q$, one has
	\[
		\mathrm{Dec}_{\F_q}(g+g')=\mathrm{Dec}_{\F_q}(g)+\mathrm{Dec}_{\F_q}(g'),
		\qquad
		\mathrm{Dec}_{\F_q}(cg)=c\,\mathrm{Dec}_{\F_q}(g),
	\]
	and $\mathrm{Dec}_{\F_q}(g_h)=h$ for every $h\in\F_q^M$.
\end{lemma}
\begin{proof}
	\leanok
	\uses{def:qld-full-field-decoder, def:low-degree-encoding,
		def:decoding-map}
	The decoder is given coordinatewise by evaluation of the representative at the points of the hypercube, and evaluation at a fixed point is $\F_q$-linear in the representative, which gives the two displayed identities. For the third, the multilinear encoding $g_h$ takes the value $h_y$ at each point $y$ of the hypercube, because the indicator polynomials of \eqref{eq:low-degree-encoding-definition} vanish at every other such point; hence the vector of its evaluations is $h$.
\end{proof}

\begin{lemma}[Evaluation identity for encoding outcomes]\label{lem:qld-decoder-evaluation}
	\lean{MIPStarRE.QPBT.decodeFq_dotProduct_indicatorVec}
	\leanok
	\uses{def:qld-full-field-decoder, def:indicator-vector}
	If $g$ is an encoding, then
	\[
		\mathrm{Dec}_{\F_q}(g)\cdot\mathrm{ind}_m(u)=g(u)
	\]
	for every $u\in\F_q^m$.
\end{lemma}
\begin{proof}
	\leanok
	\uses{def:qld-full-field-decoder, def:indicator-vector, def:low-degree-encoding}
	By \eqref{eq:low-degree-encoding-definition} the left-hand side is the value at $u$ of the multilinear encoding $g_{\mathrm{Dec}_{\F_q}(g)}$. Since $g$ is an encoding, that representative equals $g$, so the value is $g(u)$.
\end{proof}

% Local correction of the source (lines 1419 and 1492 of
% 14_analysis_of_the_pauli_basis_test.tex): the source invokes its boolean
% decoding map but describes it as returning all hypercube evaluations, and
% asserts the decoder--encoder identity for every g of individual degree at
% most d. The remark below records the corrected statements.
\begin{remark}[Choice of decoding map and the decoder--encoder identity]\label{rem:qld-decoding-identity}
	The source states the constructions of this chapter in terms of its boolean decoding map (the case $H = \{0,1\}$ of Definition~\ref{def:decoding-map}) while describing that map as returning the string of \emph{all} evaluations of $g$ on the hypercube. The proofs require the latter reading, $H = \F_q$, which is adopted throughout this chapter: the linearity of $\mathrm{Dec}$ and the identity $\mathrm{Dec}(g_h) = h$ for all $h \in \F_q^{M}$, both of which fail for the boolean map, are used in the symmetry step of the proof of Lemma~\ref{lem:qld-construct-the-paulis}. The source further asserts the identity $\mathrm{Dec}(g) \cdot \mathrm{ind}_m(u) = g(u)$ for every $g \in \mathrm{ideg}_{d,m}(\F_q)$. Since the left-hand side equals $g_{\mathrm{Dec}(g)}(u)$, the identity, taken for all $u \in \F_q^m$, states that the multilinear representative $g_{\mathrm{Dec}(g)}$ and $g$ induce the same function on $\F_q^m$. Every encoding satisfies it, because then $g_{\mathrm{Dec}(g)} = g$ as polynomial representatives; a polynomial that is not an encoding need not: for $m = 1$, $d \geq 2$, and $q > 2$, the polynomial $g(x) = x^2$ has $\mathrm{Dec}(g) \cdot \mathrm{ind}_1(u) = u$, which differs from $g(u) = u^2$ at every $u \notin \{0,1\}$. When $d \leq q-1$ the failure at some $u$ characterizes the polynomials that are not encodings, since evaluation is injective on representatives of individual degree at most $q-1$ (Definition~\ref{def:polynomials-degree}); when $d \geq q$ it does not: for the admissible tuple $(q,m,d) = (8,1,8)$, the polynomial $g(x) = x^8$ is not an encoding, yet $\mathrm{Dec}(g) \cdot \mathrm{ind}_1(u) = u = u^8 = g(u)$ for every $u \in \F_8$. Accordingly, the proofs below invoke the identity only for encodings; in the proof of Lemma~\ref{lem:qld-construct-the-paulis} the contribution of the remaining outcomes is controlled by the bound \eqref{eq:qld-nonencoding-mass} on the mass that the measurement of Lemma~\ref{lem:qld-4-7} places on polynomials that are not encodings, while in the conjugation identity \eqref{eq:qld-unitary-6} the quantity $\mathrm{Dec}(g) \cdot \mathrm{ind}_m(u)$ cancels exactly.
	The discrepancy and these two corrections are documented in
	\cite{gap:qpbt_decoding-identity}.
\end{remark}

\section{The pulled-apart measurements and observables}

The measurement $\{\hat{S}_{g_X,g_Z}\}$ reads out an $X$-polynomial and a $Z$-polynomial simultaneously. We now separate the two readings again, this time as commuting actions on distinct registers: the polynomial outcome is measured on registers $\mathrm{A}\mathrm{A}'$ (resp.\ $\mathrm{B}\mathrm{B}'$), while a genuine Pauli measurement acts on the EPR register $\mathrm{A}''$ (resp.\ $\mathrm{B}''$).

\begin{definition}[Single-basis marginals of the polynomial-pair measurement]\label{def:s-w-marginals}
	\lean{MIPStarRE.QPBT.GlobalPairWitness.marginalPoly}
	\leanok
	\uses{lem:qld-4-7}
	Let $\{\hat{S}_{g_X,g_Z}\}$ be the projective measurement of Lemma~\ref{lem:qld-4-7}. For every $g \in \mathrm{ideg}_{d,m}(\F_q)$ define
	\[
		\hat{S}^X_g = \sum_{g_Z} \hat{S}_{g,g_Z}
		\qquad \text{and} \qquad
		\hat{S}^Z_g = \sum_{g_X} \hat{S}_{g_X,g}\;,
	\]
	the sums ranging over $\mathrm{ideg}_{d,m}(\F_q)$.
\end{definition}

\begin{lemma}[Projectivity of the single-basis marginals]\label{lem:s-w-marginals-projective}
	\lean{MIPStarRE.QPBT.GlobalPairWitness.marginalPoly_isProjective}
	\leanok
	\uses{def:s-w-marginals, lem:qld-4-7}
	For each $W\in\{X,Z\}$, the family $\{\hat{S}^W_g\}_g$ is a projective measurement with outcomes in $\mathrm{ideg}_{d,m}(\F_q)$.
\end{lemma}
\begin{proof}
	\leanok
	\uses{def:s-w-marginals, lem:sandwich-postprocess-projective}
	By Definition~\ref{def:s-w-marginals}, each marginal is the finite postprocessing
	of the joint projective measurement by the outcome map $(g_X,g_Z)\mapsto g_W$.
	Lemma~\ref{lem:sandwich-postprocess-projective} shows that this postprocessing
	preserves projectivity.
\end{proof}

\begin{definition}[Pauli dot-product projectors]\label{def:tau-dot-product-projector}
	\lean{MIPStarRE.QPBT.tauDotProj}
	\leanok
	\uses{def:generalized-pauli}
	For $W \in \{X,Z\}$ and $h \in \F_q^{M}$ let $\tau^W_h$ denote the projector corresponding to measuring $(\C^q)^{\ot M}$ in the $W$ basis and obtaining the outcome $h$ (Definition~\ref{def:generalized-pauli}). For all $W \in \{X,Z\}$, $\tilde{u} \in \F_q^{M}$, and $a \in \F_q$, define the operator
	\begin{equation}\label{eq:def-tauwu}
		\tau^{W}_{a}(\tilde{u}) = \sum_{h \in \F_q^{M} \,:\, h \cdot \tilde{u} = a} \tau^W_{h}\;.
	\end{equation}
\end{definition}

\begin{lemma}[Projectivity and completeness of the dot-product operators]\label{lem:tau-dot-product-projective}
	\lean{MIPStarRE.QPBT.tauDotProj_isProj, MIPStarRE.QPBT.sum_tauDotProj_eq_one}
	\leanok
	\uses{def:tau-dot-product-projector, def:generalized-pauli}
	For each fixed $W$ and $\tilde{u}$, every $\tau^W_a(\tilde{u})$ is a projector and
	$\sum_{a\in\F_q}\tau^W_a(\tilde{u})=\Id$. Thus
	$\{\tau^W_a(\tilde{u})\}_{a\in\F_q}$ is a projective measurement.
\end{lemma}

\begin{proof}
	\leanok
	\uses{def:tau-dot-product-projector, def:generalized-pauli,
		lem:pauli-observable-expansion}
	The zero-label case of Lemma~\ref{lem:pauli-observable-expansion} has
	all phase factors equal to $1$. Since $\tau^W(0)=\Id$ by
	Definition~\ref{def:generalized-pauli}, it gives $\sum_h\tau^W_h=\Id$.
	Let $U$ be the square matrix whose columns are the $W$-basis vectors $v_h$,
	so that $\tau^W_h=v_hv_h^\dagger$. Completeness gives $UU^\dagger=\Id$,
	hence $U$ is unitary. For each $a$, let $D_a$ be the diagonal projector
	whose entry at $h$ is $1$ if $h\cdot\tilde{u}=a$ and $0$ otherwise.
	Expanding the matrix product gives
	$\tau^W_a(\tilde{u})=UD_aU^\dagger$, which is a projector.
	The fibers of $h\mapsto h\cdot\tilde{u}$ partition $\F_q^M$, so
	$\sum_a\tau^W_a(\tilde{u})=\sum_h\tau^W_h=\Id$.
\end{proof}

\begin{definition}[Pulled-apart Pauli basis measurements]\label{def:tilde-m-measurement}
	\lean{MIPStarRE.QPBT.tildeM}
	\leanok
	\uses{def:s-w-marginals, def:tau-dot-product-projector, def:qld-full-field-decoder}
	For all $W \in \{X,Z\}$ and $\tilde{u} \in \F_q^{M}$ define the measurement $\{\tilde{M}^{W,\tilde{u}}_a\}_{a \in \F_q}$ by
	\begin{equation}\label{eq:tilde_M}
		\tilde{M}^{W,\tilde{u}}_a = \sum_{g} \hat{S}^W_g \ot \tau_{(\mathrm{Dec}(g) \cdot \tilde{u}) - a}^W(\tilde{u})\;,
	\end{equation}
	the sum ranging over $g \in \mathrm{ideg}_{d,m}(\F_q)$. If $\hat{S}^W_g$ is viewed as an operator acting on registers $\mathrm{A}\mathrm{A}'$ (resp.\ $\mathrm{B}\mathrm{B}'$), then $\tilde{M}^{W,\tilde{u}}_a$ acts on registers $\mathrm{A}\mathrm{A}'\mathrm{A}''$ (resp.\ $\mathrm{B}\mathrm{B}'\mathrm{B}''$), the $\tau^W$ factor acting on $\mathrm{A}''$ (resp.\ $\mathrm{B}''$).
\end{definition}

\begin{lemma}[Projectivity and completeness of the pulled-apart measurements]\label{lem:tilde-m-projective}
	\lean{MIPStarRE.QPBT.tildeM_isProj, MIPStarRE.QPBT.sum_tildeM_eq_one}
	\leanok
	\uses{def:tilde-m-measurement}
	For each $W$ and $\tilde{u}$, every $\tilde{M}^{W,\tilde{u}}_a$ is a projector and
	$\sum_{a\in\F_q}\tilde{M}^{W,\tilde{u}}_a=\Id$. Thus
	$\{\tilde{M}^{W,\tilde{u}}_a\}_{a\in\F_q}$ is a projective measurement.
\end{lemma}

\begin{proof}
	\leanok
	\uses{def:tilde-m-measurement, lem:s-w-marginals-projective,
		lem:tau-dot-product-projective, lem:sandwich-tensor-projector,
		lem:distance-tensor-sum-right, lem:distance-tensor-sum-left}
	For fixed $W$, $\tilde{u}$, and $a$, each summand in \eqref{eq:tilde_M}
	is a tensor product of two projections, hence is a projection.
	If $g\ne g'$, then $\hat{S}^W_g\hat{S}^W_{g'}=0$ because
	$\{\hat{S}^W_g\}_g$ is a projective measurement. Thus distinct summands
	in \eqref{eq:tilde_M} are orthogonal, and their sum is a projection.
	For completeness, fix $g$ and put $c_g=\mathrm{Dec}(g)\cdot\tilde{u}$.
	The map $a\mapsto c_g-a$ is a bijection of $\F_q$, so
	\[
		\sum_a\tau^W_{c_g-a}(\tilde{u})=\Id.
	\]
	Interchanging the two finite sums in \eqref{eq:tilde_M} gives
	\[
		\sum_a\tilde{M}^{W,\tilde{u}}_a
		=\sum_g\hat{S}^W_g\ot\Id=\Id.
	\]
\end{proof}

\begin{definition}[Pulled-apart Pauli observables]\label{def:tilde-w-observables}
	\lean{MIPStarRE.QPBT.tildeObs}
	\leanok
	\uses{def:tilde-m-measurement, def:dual-self-dual-normal-basis, def:binary-representation, def:subfield-trace}
	Let $\{e_j\}_{j \in \{0,\ldots,t-1\}}$ be the self-dual normal basis for $\F_q$ over $\F_2$ (Definition~\ref{def:dual-self-dual-normal-basis}) that is fixed once and for all in Definition~\ref{def:binary-representation}, and let $\tr = \tr_{q \to 2} \colon \F_q \to \F_2$ be the subfield trace. For all $W \in \{X,Z\}$, $\tilde{u} \in \F_q^M$, and $j \in \{0,\ldots,t-1\}$ define the operator
	\begin{equation}\label{eq:def-tildewj}
		\tilde{W}^j(\tilde{u}) = \sum_{a \in \F_q} (-1)^{\tr(e_j a)} \, \tilde{M}_a^{W,\tilde{u}}\;,
	\end{equation}
	acting on the same registers as $\tilde{M}^{W,\tilde{u}}_a$.
\end{definition}

% In the source the following properties are stated inline with the
% definition of the observables; they are recorded here as a lemma.
\begin{lemma}[Product form and Pauli relations of the pulled-apart observables]\label{lem:tildew-product-form}
	\lean{MIPStarRE.QPBT.tildeObs_eq_heteroKron}
	\lean{MIPStarRE.QPBT.tildeObs_isHermitian}
	\lean{MIPStarRE.QPBT.tildeObs_mul_self}
	\lean{MIPStarRE.QPBT.tildeObs_twisted_commutation}
	\leanok
	\uses{def:tilde-w-observables, def:s-w-marginals, def:generalized-pauli, def:qld-full-field-decoder, lem:qld-4-7}
	Let $\tau^W(\cdot)$ denote the generalized Pauli $W$ observable (Definition~\ref{def:generalized-pauli}). For all $W \in \{X,Z\}$, $\tilde{u} \in \F_q^M$, and $j \in \{0,\ldots,t-1\}$,
	\begin{equation}\label{eq:tildewj-product-form}
		\tilde{W}^j(\tilde{u})
		= \Big( \sum_{g_X,g_Z} (-1)^{\tr(e_j (\mathrm{Dec}(g_W) \cdot \tilde{u}))} \, \hat{S}_{g_X,g_Z} \Big) \ot \tau^W(e_j \tilde{u})\;,
	\end{equation}
	where $g_W$ denotes $g_X$ if $W = X$ and $g_Z$ if $W = Z$. In particular each $\tilde{W}^j(\tilde{u})$ is Hermitian and squares to the identity. Moreover, the pulled-apart observables satisfy exactly the twisted commutation relations of the generalized Pauli observables $\tau^X(e_j \tilde{u})$ and $\tau^Z(e_{j'} \tilde{v})$: for all $j,j' \in \{0,\ldots,t-1\}$ and $\tilde{u},\tilde{v} \in \F_q^M$,
	\begin{equation}\label{eq:tildew-twisted-commutation}
		\tilde{X}^j(\tilde{u})\, \tilde{Z}^{j'}(\tilde{v})
		= (-1)^{\tr((e_j \tilde{u}) \cdot (e_{j'} \tilde{v}))}\, \tilde{Z}^{j'}(\tilde{v})\, \tilde{X}^j(\tilde{u})\;,
	\end{equation}
	where $\tr((e_j \tilde{u}) \cdot (e_{j'} \tilde{v})) = \tr(e_j e_{j'} (\tilde{u} \cdot \tilde{v}))$.
\end{lemma}
\begin{proof}
	\leanok
	\uses{def:tilde-w-observables, def:tilde-m-measurement, def:tau-dot-product-projector, def:s-w-marginals, lem:pauli-observable-expansion, lem:twisted-commutation, lem:qld-4-7}
	Expanding $\tau^W(e_j \tilde{u}) = \sum_{h} (-1)^{\tr(e_j (\tilde{u} \cdot h))}\, \tau^W_h$ on the right-hand side of \eqref{eq:tildewj-product-form} (Lemma~\ref{lem:pauli-observable-expansion}), grouping the $h$ according to the value $a' = h \cdot \tilde{u}$ into the projectors $\tau^W_{a'}(\tilde{u})$ of \eqref{eq:def-tauwu}, collecting the pairs $(g_X,g_Z)$ with a common value $g_W = g$ into $\hat{S}^W_g$ (Definition~\ref{def:s-w-marginals}), and changing the summation variable to $a = (\mathrm{Dec}(g) \cdot \tilde{u}) - a'$ recovers, by \eqref{eq:tilde_M} and \eqref{eq:def-tildewj}, the operator $\tilde{W}^j(\tilde{u})$; this proves \eqref{eq:tildewj-product-form}. Since the $\hat{S}_{g_X,g_Z}$ are mutually orthogonal projectors summing to the identity (Lemma~\ref{lem:qld-4-7}) and $\tau^W(e_j \tilde{u})$ is a Hermitian involution, the product form exhibits $\tilde{W}^j(\tilde{u})$ as Hermitian with square $\sum_{g_X,g_Z} \hat{S}_{g_X,g_Z} \ot \Id = \Id$. For \eqref{eq:tildew-twisted-commutation}, the two sign-weighted sums of the projectors $\hat{S}_{g_X,g_Z}$ appearing in the product forms of $\tilde{X}^j(\tilde{u})$ and $\tilde{Z}^{j'}(\tilde{v})$ commute with each other, so the commutation phase is that of $\tau^X(e_j \tilde{u})$ and $\tau^Z(e_{j'} \tilde{v})$, which is $(-1)^{\tr((e_j \tilde{u}) \cdot (e_{j'} \tilde{v}))}$ by Lemma~\ref{lem:twisted-commutation}.
\end{proof}

% Local correction of the source (lines 1451-1456 of
% 14_analysis_of_the_pauli_basis_test.tex): the source asserts commutation
% whenever j != j'.
\begin{remark}[The commutation phase across basis indices]\label{rem:qld-cross-phase}
	For $j \neq j'$ the source asserts that $\tilde{X}^j(\tilde{u})$ and $\tilde{Z}^{j'}(\tilde{v})$ commute for all $\tilde{u}, \tilde{v} \in \F_q^M$. Since $e_j e_{j'} \neq 0$ and the bilinear form $(a,b) \mapsto \tr(ab)$ on $\F_q$ is nondegenerate, there exist $\tilde{u}, \tilde{v}$ with $\tr(e_j e_{j'} (\tilde{u} \cdot \tilde{v})) = 1$, and for these \eqref{eq:tildew-twisted-commutation} yields anticommutation instead. By self-duality of the basis $\{e_j\}$, the dichotomy asserted by the source --- commutation for $j \neq j'$, and the phase $(-1)^{\tr((e_j\tilde{u}) \cdot (e_j\tilde{v}))}$ for $j = j'$ --- agrees with \eqref{eq:tildew-twisted-commutation} whenever $\tilde{u} \cdot \tilde{v} \in \{0,1\}$, in particular for all $\tilde{u}, \tilde{v} \in \{0,1\}^M$. The relations \eqref{eq:tildew-twisted-commutation} are not used in the remainder of the chapter.
	The counterexample and corrected phase are documented in
	\cite{gap:qpbt_cross-basis-phase}.
\end{remark}

\section{Consistency with the points measurements}

The next lemma shows that the pulled-apart measurements are consistent with the original points measurements of the strategy, and self-consistent. Write $\delta_G$ for the error of the global polynomial-pair witness supplied by Lemma~\ref{lem:qld-4-7}. Assume $0\leq\eps\leq1$ and $\delta_G\geq0$, and put
\[
	B = \delta_G + \sqrt{\eps} + \frac{md}{q}.
\]
Each displayed construction estimate below has its own universal constant $C_E \geq 1$, quantified independently, and its own scale $\delta_P^E=C_EB$. The source uses one enlarged error $\delta_S$ for all of these estimates. The two presentations agree: after the finitely many estimates have been established, take $C_* = \max_E C_E$ and $\delta_P=C_*B$. Since $B\geq0$, every bound remains valid at this common scale. In the proof paragraphs, $\delta_P$ denotes this common enlargement. It has the same asymptotic form as $\delta_G$: after shrinking $b$ to at most $1/2$, the term $\eps^b$ dominates $\sqrt{\eps}$ for $0 \leq \eps \leq 1$, while $(md)^a q^{-b}$ dominates $md/q$.

\begin{lemma}[Consistency and self-consistency of the pulled-apart measurements]\label{lem:qld-construct-the-paulis}
	\lean{MIPStarRE.QPBT.exists_pulled_apart_consistency}
	\lean{MIPStarRE.QPBT.deltaConstructPaulis}
	\leanok
	\uses{def:tilde-m-measurement, def:tilde-w-observables,
		lem:projective-strategy-setup, def:expanded-state, def:consistency,
		def:approx-question-indexed-operators, lem:qld-4-7,
		def:indicator-vector}
	There are universal constants $a>1$, $0<b<1$, and $C_*\geq1$ such that,
	for every admissible parameter tuple and projective strategy in the setting
	above with $0\leq\eps\leq1$, there is a pair of global projective measurements
	satisfying Lemma~\ref{lem:qld-4-7} at error $\delta_G(\eps,m,d,q)$.
	The measurements $\{\tilde{M}^{W,\tilde{u}}_a\}_{a \in \F_q}$ and observables
	$\tilde{W}^j(\tilde{u})$ constructed from this same pair satisfy the following
	properties, with all three constants taken to be $C_*$.
	% The paper's statement also quantifies item 1 over j, but j does not occur
	% in item 1; the vacuous quantifier is dropped here.
	\begin{enumerate}
		\item (\textbf{Consistency with points measurements}) For all $W \in \{X,Z\}$, on average over uniformly random $u \in \F_q^m$,
		\[
			\Id_{\mathrm{A}\mathrm{A}'\mathrm{A}''} \ot \tilde{M}^{W,\mathrm{ind}_m(u)}_a \simeq_{\delta_P^{\mathrm{AB}}} M^{(\mathrm{Point},W),u}_a \ot \Id_{\mathrm{B}\mathrm{B}'\mathrm{B}''}
		\]
		and
		\[
			\tilde{M}^{W,\mathrm{ind}_m(u)}_a \ot \Id_{\mathrm{B}\mathrm{B}'\mathrm{B}''} \simeq_{\delta_P^{\mathrm{BA}}} \Id_{\mathrm{A}\mathrm{A}'\mathrm{A}''} \ot M^{(\mathrm{Point},W),u}_a\;,
		\]
		with the answer summation over $a \in \F_q$.
		\item (\textbf{Self-consistency}) For all $W \in \{X,Z\}$ and $j \in \{0,\ldots,t-1\}$, on average over uniformly random $\tilde{u} \in \F_q^{M}$,
		\[
			\tilde{W}^j(\tilde{u}) \ot \Id_{\mathrm{B}\mathrm{B}'\mathrm{B}''} \approx_{\delta_P^{\mathrm{obs}}} \Id_{\mathrm{A}\mathrm{A}'\mathrm{A}''} \ot \tilde{W}^j(\tilde{u})\;.
		\]
	\end{enumerate}
	The construction above does not assume a supplied global measurement. Its
	proof chooses the global polynomial-pair measurements supplied by
	Lemma~\ref{lem:qld-4-7}, applies the three estimates of
	Example~\ref{lem:qld-construct-the-paulis-given-global-pair} to the same
	witness, and replaces their universal constants by their maximum. The
	separate numerical and isometry-transfer questions arising later in the
	extraction argument are recorded in \cite{gap:qpbt_extraction-transfer}.
\end{lemma}
\begin{proof}
	\leanok
	\uses{lem:qld-4-7, lem:qld-construct-the-paulis-given-global-pair}
	Choose the global polynomial-pair measurements from
	Lemma~\ref{lem:qld-4-7}. Example~\ref{lem:qld-construct-the-paulis-given-global-pair}
	provides universal constants for the two point-consistency estimates and the
	observable self-consistency estimate for this same witness. Let $C_*$ be their
	maximum. Since $B=\delta_G+\sqrt\eps+md/q$ is nonnegative, replacing each
	constant by $C_*$ only enlarges its allowed error. The three conclusions
	therefore hold simultaneously at the common scale $C_*B$.
\end{proof}

\begin{example}[Conditional estimates given a global polynomial-pair witness]
	\label{lem:qld-construct-the-paulis-given-global-pair}
	\lean{MIPStarRE.QPBT.tildeM_consistent_pointMeas_ofGlobalPairWitness}
	\lean{MIPStarRE.QPBT.tildeM_consistent_pointMeas'_ofGlobalPairWitness}
	\lean{MIPStarRE.QPBT.tildeObs_selfConsistent_ofGlobalPairWitness}
	\leanok
	\uses{def:tilde-m-measurement, def:tilde-w-observables, lem:qld-4-7,
		lem:projective-strategy-setup, def:expanded-state, def:consistency,
		def:approx-question-indexed-operators, def:indicator-vector}
	For each of the three estimates in Lemma~\ref{lem:qld-construct-the-paulis},
	there is a universal constant $C_E\geq1$ such that the estimate holds for
	every admissible tuple, every $0\leq\eps\leq1$, every projective strategy at
	error $\eps$, every $\delta_G\geq0$, and every supplied pair of projective
	polynomial-pair measurements satisfying both point-consistency estimates of
	Lemma~\ref{lem:qld-4-7} at error $\delta_G$.
	The global measurement and its consistency estimates are hypotheses of this
	statement. Consequently, this auxiliary does not assert their construction
	or the unrestricted conclusion of Lemma~\ref{lem:qld-construct-the-paulis}.
\end{example}
\begin{proof}
	\uses{lem:qld-4-7, lem:qld-constructing-the-paulis-helper, def:expanded-point-measurement, def:tilde-m-measurement, def:tau-dot-product-projector, lem:qld-decoder-evaluation, lem:qld-decoder-linearity, def:indicator-vector, lem:qld-win-implications, lem:qld-comm-cons, fact:add-a-proj, fact:agreement, fact:triangle, fact:data-processing, def:bracket, lem:povm-to-obs, lem:schwartz-zippel-total-degree}
	% Support restriction added: the source (line 1492 of
	% 14_analysis_of_the_pauli_basis_test.tex) rewrites Dec(g).ind_m(u) as
	% g(u) for every outcome g; see Remark rem:qld-decoding-identity.
	We first bound the mass that the marginal measurements $\{\hat{S}^W_g\}$ place on polynomials that are not encodings. On average over uniformly random $u \in \F_q^m$, the point-consistency relation \eqref{eq:qld-s-point-con-bob} of Lemma~\ref{lem:qld-4-7} gives $(\hat{S}^W_{[g(u)=a]})_{\mathrm{B}\mathrm{B}'} \approx_{\delta_G} (\hat{M}^{(\mathrm{Point},W),u}_a)_{\mathrm{A}\mathrm{B}''}$ (Lemma~\ref{fact:agreement}); the consistency and Pauli-basis-consistency checks of Lemma~\ref{lem:qld-win-implications} give, again on average over uniform $u$ and combining through Lemma~\ref{fact:data-processing}, Lemma~\ref{fact:agreement}, and Lemma~\ref{fact:triangle}, that $(M^{(\mathrm{Point},W),u}_{a'})_{\mathrm{A}} \approx_{O(\eps)} (M^{(\mathrm{Pauli},W)}_{[g_h(u)=a']})_{\mathrm{A}}$; and replacing the point factor of the convolution form of $\hat{M}^{(\mathrm{Point},W),u}_a$ (Definition~\ref{def:expanded-point-measurement}) accordingly costs the same error, because the projectors $\tau^W_{a''}(\mathrm{ind}_m(u))$ of \eqref{eq:def-tauwu} are mutually orthogonal (Lemma~\ref{fact:add-a-proj}). Writing
	\[
		K^u_a = \sum_{a' + a'' = a} \bigl(M^{(\mathrm{Pauli},W)}_{[g_h(u)=a']}\bigr)_{\mathrm{A}} \ot \bigl(\tau^W_{a''}(\mathrm{ind}_m(u))\bigr)_{\mathrm{B}''}\;,
	\]
	a projective measurement, Lemma~\ref{fact:triangle} and Lemma~\ref{fact:agreement} therefore give
	\[
		\E_{u \in \F_q^m} \sum_a \bra{\hat\psi} (\hat{S}^W_{[g(u)=a]})_{\mathrm{B}\mathrm{B}'} \ot (K^u_a)_{\mathrm{A}\mathrm{B}''} \ket{\hat\psi} \geq 1 - O(\delta_G + \eps)\;.
	\]
	Expanding $K^u_a$ over pairs $(h, h'')$ of Pauli outcomes, a term of the left-hand side with outcomes $(g, h, h'')$ is constrained by $g(u) = g_h(u) + g_{h''}(u) = g_{h + h''}(u)$, by the linearity of the encoding \eqref{eq:low-degree-encoding-definition}. If $g$ is not an encoding then $g \neq g_{h+h''}$, and both are polynomials of total degree at most $md$, so the constraint holds for at most an $md/q$ fraction of $u$ (Lemma~\ref{lem:schwartz-zippel-total-degree}). Since all terms are nonnegative and, for each $g$, the operators of the $(h,h'')$-sum add up to the identity, it follows that
	\begin{equation}\label{eq:qld-nonencoding-mass}
	\sum_{g \,\text{not an encoding}} \bra{\hat\psi} (\hat{S}^W_g)_{\mathrm{B}\mathrm{B}'} \ket{\hat\psi} \,\leq\, O(\delta_G + \eps) + \frac{md}{q} \,\leq\, \delta_P\;,
	\end{equation}
	and symmetrically with the registers $\mathrm{A}\mathrm{A}'$ in place of $\mathrm{B}\mathrm{B}'$. (If $d = 0$, every element of $\mathrm{ideg}_{d,m}(\F_q)$ is constant, hence an encoding, and the left-hand side vanishes.)

	For the first item, in the case where $M^{(\mathrm{Point},W),u}_a$ acts on register $\mathrm{A}$ and $\tilde{M}^{W,\mathrm{ind}_m(u)}_a$ on $\mathrm{B}\mathrm{B}'\mathrm{B}''$, expanding \eqref{eq:tilde_M} and \eqref{eq:def-tauwu} and using the convolution form of the combined projectors $\hat{M}^{(\mathrm{Point},W),u}_a$ (Definition~\ref{def:expanded-point-measurement}) gives the exact identity
	\[
		\E_{u \in \F_q^m} \sum_a \bra{\hat\psi} M^{(\mathrm{Point},W),u}_a \ot \tilde{M}^{W,\mathrm{ind}_m(u)}_a \ket{\hat\psi}
		= \E_{u \in \F_q^m} \sum_g \bra{\hat\psi} (\hat{M}^{(\mathrm{Point},W),u}_{\mathrm{Dec}(g) \cdot \mathrm{ind}_m(u)})_{\mathrm{A}\mathrm{B}''} \ot (\hat{S}^W_g)_{\mathrm{B}\mathrm{B}'} \ket{\hat\psi}\;.
	\]
	For encodings $g$ the subscript $\mathrm{Dec}(g) \cdot \mathrm{ind}_m(u)$ equals $g(u)$; discarding the (nonnegative) terms with $g$ not an encoding and reinstating the corresponding terms with subscript $g(u)$, each of which is at most $\bra{\hat\psi} (\hat{S}^W_g)_{\mathrm{B}\mathrm{B}'} \ket{\hat\psi}$, shows that the right-hand side is at least
	\[
		\E_{u \in \F_q^m} \sum_g \bra{\hat\psi} (\hat{M}^{(\mathrm{Point},W),u}_{g(u)})_{\mathrm{A}\mathrm{B}''} \ot (\hat{S}^W_g)_{\mathrm{B}\mathrm{B}'} \ket{\hat\psi} - \delta_P
		\geq 1 - \delta_P
	\]
	by \eqref{eq:qld-nonencoding-mass} and Lemma~\ref{lem:qld-4-7}, and the register-interchanged relation follows analogously. For the second item one first shows that, on average over uniformly random $\tilde{u} \in \F_q^M$,
	\begin{equation}\label{eq:qld-pulling-cons}
		\bigl( \tilde{M}^{W,\tilde{u}}_a \bigr)_{\mathrm{A}\mathrm{A}'\mathrm{A}''} \approx_{\delta_P} \bigl( \tilde{M}^{W,\tilde{u}}_a \bigr)_{\mathrm{B}\mathrm{B}'\mathrm{B}''}\;:
	\end{equation}
	right-multiplying by the resolution $\sum_g (\hat{S}^W_g)_{\mathrm{A}\mathrm{A}'} \ot (\hat{M}^{(\mathrm{Point},W),u}_{g(u)})_{\mathrm{B}\mathrm{A}''} \approx_{\delta_P} \Id$ of Lemma~\ref{lem:qld-constructing-the-paulis-helper} (with Lemma~\ref{fact:add-a-proj}) and expanding the products of $\tau^W$ projectors via \eqref{eq:def-tauwu} and $h \cdot \mathrm{ind}_m(u) = g_h(u)$ places $(\tilde{M}^{W,\tilde{u}}_a)_{\mathrm{A}\mathrm{A}'\mathrm{A}''}$ within $O(\delta_P)$ of a sum over pairs $(g,h)$ with $(\mathrm{Dec}(g)-h)\cdot\tilde{u} = a$. One then applies, in turn, the self-consistency of the points measurements (Lemma~\ref{lem:qld-win-implications}) at cost $\eps$; the exact identity $\ket{\hat\psi} = \sum_{h'} (\tau^W_{h'})_{\mathrm{B}'} \ot (\tau^W_{h'})_{\mathrm{B}''} \ket{\hat\psi}$ of the EPR registers; Lemma~\ref{lem:qld-constructing-the-paulis-helper} twice more (once through a Cauchy--Schwarz detour of cost $O(\sqrt{\eps})$ using the self-consistency of the $\hat{M}^{(\mathrm{Point},W),u}_a$, Lemma~\ref{lem:qld-comm-cons}); and finally a Schwartz--Zippel step (Lemma~\ref{lem:schwartz-zippel-total-degree}, for polynomials of total degree at most $md$) which replaces the event $(g'-g_{h'})(u) = (g-g_h)(u)$ by the polynomial identity $g'-g_{h'} = g-g_h$ at cost $md/q$. The resulting expression
	\[
		\sum_{\substack{g,h,g',h': \\ g' - g_{h'} = g - g_h \\ (\mathrm{Dec}(g)-h)\cdot \tilde{u} = a}}
		\Bigl((\hat{S}^W_g)_{\mathrm{A}\mathrm{A}'} \ot (\tau_h^W)_{\mathrm{A}''} \Bigr) \cdot \Bigl( (\hat{S}^W_{g'})_{\mathrm{B}\mathrm{B}'} \ot (\tau_{h'}^W)_{\mathrm{B}''} \Bigr)
	\]
	is symmetric under exchanging $(g,h)$ with $(g',h')$: since $\mathrm{Dec}$ is linear with $\mathrm{Dec}(g_{h}) = h$ for all $h \in \F_q^M$ (Remark~\ref{rem:qld-decoding-identity}), the constraint $g' - g_{h'} = g - g_h$ implies $\mathrm{Dec}(g) - h = \mathrm{Dec}(g') - h'$, so the two constraint sets coincide. The same derivation therefore bounds the distance of the displayed expression to $(\tilde{M}^{W,\tilde{u}}_a)_{\mathrm{B}\mathrm{B}'\mathrm{B}''}$, and \eqref{eq:qld-pulling-cons} follows at the common scale $\delta_P$. Finally, since $\{\tilde{M}^{W,\tilde{u}}_a\}$ is projective, Item~2 of Lemma~\ref{fact:agreement}, data processing along the map $a \mapsto \tr(e_j a)$ (Lemma~\ref{fact:data-processing}), and Item~1 of Lemma~\ref{fact:agreement} yield $(\tilde{M}^{W,\tilde{u}}_{[\tr(e_j \cdot) = b]})_{\mathrm{A}\mathrm{A}'\mathrm{A}''} \approx_{\delta_P} (\tilde{M}^{W,\tilde{u}}_{[\tr(e_j \cdot) = b]})_{\mathrm{B}\mathrm{B}'\mathrm{B}''}$ with answer summation over $b \in \F_2$, and Lemma~\ref{lem:povm-to-obs} converts this into the claimed self-consistency of the observables $\tilde{W}^j(\tilde{u})$.
\end{proof}

% Source: references/qpbt-paper/14_analysis_of_the_pauli_basis_test.tex:1594--1602.
% The same scalar estimate is used before eq:qld-nonencoding-mass.
% tracked in issue #259
\begin{lemma}[Formalization support for polynomial evaluation collisions]
\label{lem:qld-poly-eval-collision}
	\lean{MIPStarRE.QPBT.poly_eval_collision_le}
	\leanok
	\uses{def:admissible, def:polynomials-degree}
	This auxiliary lemma isolates the scalar Schwartz--Zippel estimate used in
	the proof of Lemma~\ref{lem:qld-construct-the-paulis}
	in~\cite{Ji2020MIPStar}; it is not a separately named result of the source.
	Let $(q,m,d)$ be admissible parameters and let
	$g,g'\in\mathrm{ideg}_{d,m}(\F_q)$ be distinct polynomial representatives.
	For uniformly distributed $u\in\F_q^m$,
	\[
		\Prb_{u\sim\F_q^m}[g(u)=g'(u)]
		=\E_{u\in\F_q^m}\mathbf{1}[g(u)=g'(u)]
		\leq\frac{md}{q}.
	\]
\end{lemma}
\begin{proof}
	\leanok
	\uses{lem:schwartz-zippel-individual}
	Apply Lemma~\ref{lem:schwartz-zippel-individual} to $g$ and $g'$.
	The uniform average of the agreement indicator is the fraction of points
	at which the two evaluations agree, and the field has cardinality $q$.
\end{proof}

\begin{lemma}[Non-encoding support estimate for a supplied witness]
\label{lem:qld-nonencoding-mass-bound}
	\lean{MIPStarRE.QPBT.nonencodingMarginalMass}
	\lean{MIPStarRE.QPBT.nonencodingMarginalMass_le}
	\leanok
	\uses{def:s-w-marginals, def:qld-full-field-decoder, def:expanded-state,
		lem:qld-4-7}
	This formalization-only auxiliary estimate is not a separate statement of
	the source paper. Fix a global polynomial-pair witness with projective
	measurements and both point-consistency bounds of Lemma~\ref{lem:qld-4-7},
	with error $\delta_G$. If $0\leq\eps\leq1$ and $\delta_G\geq0$, then for
	either player and either Pauli basis, the non-encoding marginal mass in
	Equation~\eqref{eq:qld-nonencoding-mass} is at most
	$\delta_P^{\mathrm{supp}}=C_{\mathrm{supp}}B$ for a universal
	$C_{\mathrm{supp}}\geq1$, independent of the parameters and the witness.
	The estimate is not an additional hypothesis of the extraction lemma.
\end{lemma}
\begin{proof}
	\leanok
	\uses{lem:qld-win-implications, lem:qld-comm-cons,
		fact:triangle-for-simeq, lem:schwartz-zippel-individual}
	Measure the original Pauli answer $h$ and the ancillary Pauli answer $h'$
	and return $g_{h+h'}$. This reference measurement is supported on encodings.
	Perfect ancillary consistency preserves the point and Pauli-basis
	consistency defects under additive convolution. The consistency triangle
	therefore bounds the evaluated inconsistency between the supplied marginal
	and the reference by $\delta_G+c\sqrt\eps$, for a universal $c$.
	Schwartz--Zippel bounds their full polynomial inconsistency by
	$\delta_G+c\sqrt\eps+md/q$. Every non-encoding marginal outcome differs
	from the reference outcome, so its total mass satisfies this bound.
	Enlarging the universal constant gives $C_{\mathrm{supp}}B$ on both sides.
	The global witness is an input to this auxiliary estimate; the proof of
	Lemma~\ref{lem:qld-construct-the-paulis} constructs its own witness and
	uses separate point-consistency estimates.
\end{proof}

% Following the source, the next lemma is stated after
% lem:qld-construct-the-paulis, whose proof invokes it; its own proof
% depends only on lem:qld-4-7 and lem:qld-comm-cons, so there is no
% circularity.
Lemmas~\ref{lem:qld-4-7} and~\ref{lem:qld-comm-cons} imply the following two estimates, which were used in the proof of Lemma~\ref{lem:qld-construct-the-paulis}.

\begin{lemma}[The marginals absorb the point measurements]\label{lem:qld-constructing-the-paulis-helper}
	\lean{MIPStarRE.QPBT.sum_marginalPoly_pointMeas_approx_id}
	\lean{MIPStarRE.QPBT.marginalPoly_sub_pointMeas_approx_zero}
	\leanok
	\uses{def:s-w-marginals, def:expanded-point-measurement, def:expanded-state, def:povm-distance, lem:qld-4-7, def:symmetric-equivalents}
	For $W \in \{X,Z\}$, on average over uniformly random $u \in \F_q^m$,
	\begin{equation}\label{eq:qld-sg-cons}
		\sum_g (\hat{S}^W_g)_{\mathrm{A}\mathrm{A}'} \ot (\hat{M}^{(\mathrm{Point},W),u}_{g(u)})_{\mathrm{B}\mathrm{A}''} \approx_{\delta_P^{\mathrm{id}}} \Id
	\end{equation}
	and
	\begin{equation}\label{eq:qld-sg-cons2}
		\bigl(\hat{S}^W_g \cdot (\Id - \hat{M}^{(\mathrm{Point}, W), u}_{g(u)})\bigr)_{\mathrm{A}\mathrm{A}'} \approx_{\delta_P^{\mathrm{ann}}} 0\;,
	\end{equation}
	where the answer summation is over $g \in \mathrm{ideg}_{d,m}(\F_q)$. The independently quantified scale $\delta_P^{\mathrm{ann}}$ includes the additional point-measurement error in \eqref{eq:qld-sg-cons2}. Furthermore, all symmetric equivalents of these approximations also hold (Definition~\ref{def:symmetric-equivalents}).
\end{lemma}
\begin{proof}
	\uses{lem:qld-4-7, fact:agreement, fact:add-a-proj, def:bracket, lem:qld-comm-cons}
	The point-consistency conclusion of Lemma~\ref{lem:qld-4-7} states that, on average over $u \in \F_q^m$,
	\[
		\bra{\hat\psi} \sum_g (\hat{S}^W_g)_{\mathrm{A}\mathrm{A}'} \ot (\hat{M}^{(\mathrm{Point}, W), u}_{g(u)})_{\mathrm{B}\mathrm{A}''} \ket{\hat\psi} \geq 1 - \delta_G\;,
	\]
	which says exactly that the operator on the left-hand side of \eqref{eq:qld-sg-cons} is $\delta_G$-consistent with $\Id$; Lemma~\ref{fact:agreement} and the universal-factor enlargement give the $\approx_{\delta_P}$ statement \eqref{eq:qld-sg-cons}. For \eqref{eq:qld-sg-cons2}, the same consistency, written with the coarse-graining $\hat{S}^W_{[g(u)=a]} = \sum_{g : g(u) = a} \hat{S}^W_g$ as $(\hat{S}^W_{[g(u)=a]})_{\mathrm{A}\mathrm{A}'} \approx_{O(\delta_G)} (\hat{M}^{(\mathrm{Point}, W), u}_{a})_{\mathrm{B}\mathrm{A}''}$ with answer summation over $a \in \F_q$, is left-multiplied by $(\hat{S}^W_{[g(u)=a]})_{\mathrm{A}\mathrm{A}'}$ using Lemma~\ref{fact:add-a-proj}; moving $\hat{M}^{(\mathrm{Point}, W), u}_a$ from registers $\mathrm{B}\mathrm{A}''$ to $\mathrm{A}\mathrm{A}'$ costs a further $\eps$ by Cauchy--Schwarz together with the self-consistency of the $\hat{M}^{(\mathrm{Point}, W), u}_a$ (Lemma~\ref{lem:qld-comm-cons}), giving
	\[
		(\hat{S}^W_{[g(u)=a]})_{\mathrm{A}\mathrm{A}'} \approx_{O(\delta_G + \eps)} (\hat{S}^W_{[g(u)=a]} \cdot \hat{M}^{(\mathrm{Point}, W), u}_{a})_{\mathrm{A}\mathrm{A}'}\;.
	\]
	Expanding
	\[
		\E_u \sum_g \bigl\| (\hat{S}^W_{g} \cdot (\Id - \hat{M}^{(\mathrm{Point}, W), u}_{g(u)}))_{\mathrm{A}\mathrm{A}'} \ket{\hat\psi} \bigr\|^2
		= \E_u \sum_a \bra{\hat\psi} \bigl((\Id - \hat{M}^{(\mathrm{Point}, W), u}_{a}) \cdot \hat{S}^W_{[g(u)=a]} \cdot (\Id - \hat{M}^{(\mathrm{Point}, W), u}_{a})\bigr)_{\mathrm{A}\mathrm{A}'} \ket{\hat\psi}
	\]
	and bounding the right-hand side by $O(\delta_G + \eps)$ using the two approximations above yields \eqref{eq:qld-sg-cons2} at scale $\delta_P$. The symmetric equivalents follow by the same argument with the register bipartitions interchanged.
\end{proof}

	\section{Formalization support for the encoding-supported reference measurement}

	% Internal auxiliary constructions for eq:qld-nonencoding-mass. They are not
	% the quantitative support estimate of lem:qld-nonencoding-mass-bound, and
	% they neither use a winning premise nor construct a global polynomial-pair
	% measurement.

	\begin{definition}[The ideal Pauli register measurement]
		\label{def:qld-pauli-register-measurement}
		\lean{MIPStarRE.QPBT.pauliRegisterMeas}
		\leanok
		This measurement is not named in the source article. For a fixed Pauli
		basis it is the measurement of the ancillary Pauli register in that basis,
		assembled from the positivity and completeness of the Pauli projectors.
	\end{definition}

	\begin{lemma}[Bounded encodings lie in the decoder's image]
		\label{lem:qld-is-encoding-encoding-poly}
		\lean{MIPStarRE.QPBT.isEncoding_encodingPoly}
		\leanok
		\uses{def:low-degree-encoding, def:qld-full-field-decoder}
		The low-degree encoding of Definition~\ref{def:low-degree-encoding} of any
		Pauli register value lies in the encoding image of the full-field decoder
		of Definition~\ref{def:qld-full-field-decoder}.
	\end{lemma}
	\begin{proof}
		\leanok
		\uses{lem:qld-decoder-linearity}
		Decoding the encoding of a register value returns that value.
	\end{proof}

	\begin{definition}[The encoding-supported Pauli reference measurement]
		\label{def:qld-encoding-pauli-measurement}
		\lean{MIPStarRE.QPBT.ProjectiveSetting.encodingPauliMeas}
		\leanok
		\uses{def:qld-pauli-register-measurement, def:low-degree-encoding,
			def:post-processing}
		This measurement is not named in the source article; it is the internal
		reference measurement underlying Equation~\eqref{eq:qld-nonencoding-mass}.
		For a player and a Pauli basis, measure the strategy's Pauli answer
		together with the ideal Pauli answer of
		Definition~\ref{def:qld-pauli-register-measurement} on the ancillary
		register, and return the low-degree encoding of the sum of the two answers.
		This is not the global polynomial-pair measurement whose existence
		Lemma~\ref{lem:qld-4-7} asserts.
	\end{definition}

	\begin{lemma}[The reference effects are sums over encoding fibers]
		\label{lem:qld-encoding-pauli-effect}
		\lean{MIPStarRE.QPBT.ProjectiveSetting.encodingPauliMeas_effect}
		\leanok
		\uses{def:qld-encoding-pauli-measurement}
		The effect of the measurement of
		Definition~\ref{def:qld-encoding-pauli-measurement} at a polynomial is the
		sum of the product effects over the pairs of Pauli answers whose sum
		encodes that polynomial.
	\end{lemma}
	\begin{proof}
		\leanok
		This is the definition of the relabeled product measurement.
	\end{proof}

	\begin{lemma}[The reference measurement is a complete POVM]
		\label{lem:qld-encoding-pauli-complete}
		\lean{MIPStarRE.QPBT.ProjectiveSetting.encodingPauliMeas_effect_nonneg}
		\lean{MIPStarRE.QPBT.ProjectiveSetting.sum_encodingPauliMeas_effect_eq_one}
		\leanok
		\uses{def:qld-encoding-pauli-measurement}
		Every effect of the measurement of
		Definition~\ref{def:qld-encoding-pauli-measurement} is positive
		semidefinite, and the effects sum to the identity on the strategy and
		ancillary registers.
	\end{lemma}
	\begin{proof}
		\leanok
		Positivity and completeness are inherited from the product measurement and
		are preserved by relabeling of the outcomes.
	\end{proof}

	\begin{lemma}[Exact encoding support of the reference measurement]
		\label{lem:qld-encoding-pauli-support}
		\lean{MIPStarRE.QPBT.ProjectiveSetting.encodingPauliMeas_effect_eq_zero_of_not_isEncoding}
		\leanok
		\uses{def:qld-encoding-pauli-measurement, def:qld-full-field-decoder}
		The effect of the measurement of
		Definition~\ref{def:qld-encoding-pauli-measurement} vanishes at every
		polynomial outside the encoding image of the full-field decoder of
		Definition~\ref{def:qld-full-field-decoder}. This is exact support of the
		auxiliary measurement; it is not a bound on the non-encoding marginal mass
		of Lemma~\ref{lem:qld-nonencoding-mass-bound}.
	\end{lemma}
	\begin{proof}
		\leanok
		\uses{lem:qld-encoding-pauli-effect, lem:qld-is-encoding-encoding-poly}
		By Lemma~\ref{lem:qld-encoding-pauli-effect} the effect is a sum over the
		pairs of Pauli answers whose sum encodes the given polynomial. By
		Lemma~\ref{lem:qld-is-encoding-encoding-poly} every such polynomial lies in
		the encoding image, so outside that image the index set is empty.
	\end{proof}

	\begin{lemma}[Formalization support for mass outside a supported reference measurement]
		\label{lem:qld-supported-reference-point-mass}
		\lean{MIPStarRE.QPBT.mass_outside_support_le_point_defect}
		\leanok
		\uses{def:povm-conventions, def:distance-state-quadratic-form}
		This is not a named statement of \cite{Ji2020MIPStar}; it isolates the
		support comparison used in Equation~\eqref{eq:qld-nonencoding-mass}. Let
		$M=\{M_a\}_{a\in\mathcal A}$ and $R=\{R_a\}_{a\in\mathcal A}$ be
		complete measurements on two finite-dimensional spaces, let
		$E\subseteq\mathcal A$, and suppose $R_a=0$ for every $a\notin E$.
		For every vector $\ket\psi$ in the tensor product, without a normalization
		assumption,
		\[
			\sum_{a\notin E}q_\psi(M_a\ot\Id)
			\leq \sum_{a\ne b}q_\psi(M_a\ot R_b).
		\]
	\end{lemma}
	\begin{proof}
		\leanok
		Use completeness of $R$ to expand each term on the left as a sum over its
		outcomes. When $a\notin E$, the diagonal term vanishes by the support
		hypothesis. Every remaining quadratic form is nonnegative, so restricting
		the first outcome to $a\notin E$ can only decrease the full off-diagonal
		sum.
	\end{proof}

	\begin{lemma}[Formalization support for averaged mass outside a supported reference measurement]
		\label{lem:qld-supported-reference-average-mass}
		\lean{MIPStarRE.QPBT.avg_mass_outside_support_le_consistency_defect}
		\leanok
		\uses{def:consistency, def:distance-state-quadratic-form}
		Let $\mu$ be any finite-support nonnegative weighting on a finite question
		set, not necessarily of total mass one. For each question $x$, let $M^x$
		and $R^x$ be complete measurements with common finite outcome set
		$\mathcal A$, and suppose $(R^x)_a=0$ for every $a\notin E$. Then, for
		every vector $\ket\psi$ and without a normalization assumption,
		\[
			\operatorname{Avg}_{x\sim\mu}
			\left(\sum_{a\notin E}q_\psi(M^x_a\ot\Id)\right)
			\leq
			\operatorname{Avg}_{x\sim\mu}
			\left(\sum_{a\ne b}q_\psi(M^x_a\ot R^x_b)\right).
		\]
		The right-hand side is the bipartite consistency defect of the two
		measurement families.
	\end{lemma}
	\begin{proof}
		\leanok
		\uses{lem:qld-supported-reference-point-mass}
		Apply Lemma~\ref{lem:qld-supported-reference-point-mass} at each question
		and average the resulting inequalities with the nonnegative weights of
		$\mu$.
	\end{proof}

	\begin{lemma}[Evaluation splits over the two Pauli answers]
		\label{lem:qld-encoding-pauli-eval}
		\lean{MIPStarRE.QPBT.ProjectiveSetting.encodingPauliMeas_postprocess_eval}
		\leanok
		\uses{def:qld-encoding-pauli-measurement, def:low-degree-encoding,
			def:post-processing}
		Evaluating the polynomial outcome of
		Definition~\ref{def:qld-encoding-pauli-measurement} at a point is the same
		as evaluating the two Pauli answers separately at that point and adding the
		two values.
	\end{lemma}
	\begin{proof}
		\leanok
		\uses{lem:qld-tensor-measurement-postprocess}
		Composing the two relabelings and moving them through the product
		measurement with Lemma~\ref{lem:qld-tensor-measurement-postprocess} reduces
		the claim to additivity of the low-degree encoding of
		Definition~\ref{def:low-degree-encoding} in its register argument, which
		holds because that encoding is a scalar product with the point.
	\end{proof}

	\begin{lemma}[The evaluated reference effects are the Pauli-answer convolution]
		\label{lem:qld-encoding-pauli-convolution}
		\lean{MIPStarRE.QPBT.ProjectiveSetting.encodingPauliMeas_eval_effect_eq_convolution}
		\leanok
		\uses{def:qld-encoding-pauli-measurement, def:expanded-point-measurement}
		The effect at a field value of the evaluated measurement of
		Definition~\ref{def:qld-encoding-pauli-measurement} is the sum, over the
		pairs of field values adding to that value, of the products of the
		evaluated strategy effect and the ideal point projector. This is the
		convolution used in the proof of
		Equation~\eqref{eq:qld-nonencoding-mass}; it asserts neither an
		approximation nor consistency with the strategy's point measurement.
	\end{lemma}
	\begin{proof}
		\leanok
		\uses{lem:qld-encoding-pauli-eval}
		Expand the relabeled product measurement of
		Lemma~\ref{lem:qld-encoding-pauli-eval} at the given field value; the
		ancillary factor is the ideal point projector by definition of the
		low-degree encoding.
	\end{proof}

\section{The swap unitaries and the extraction of the Pauli group}

The swap unitaries are controlled operators: a complete projective
measurement on the polynomial registers selects, outcome by outcome, a product
of Pauli observables on the register $\mathrm{A}''$ (resp.\ $\mathrm{B}''$).
The following algebraic statement is not a named result of the source article;
it isolates the orthogonality calculation that the source performs inline, and
the proof of Lemma~\ref{lem:v-swap-conjugation} cites it for both the unitarity
of $V_{\mathrm{A}}$ and the outcomewise form of the conjugation.

\begin{lemma}[Formalization support for controlled operator sums]
	\label{lem:controlled-sum-support}
	\lean{MIPStarRE.Quantum.sum_heteroKron_mul_sum_heteroKron}
	\lean{MIPStarRE.Quantum.sum_heteroKron_mul_conjTranspose}
	\leanok
	\uses{def:povm-conventions}
	Let $\{S_o\}_{o \in \mathcal{O}}$ be a projective measurement on a
	finite-dimensional space, so that the $S_o$ are pairwise orthogonal
	projectors with $\sum_o S_o = \Id$, and let $A_o$ and $B_o$ be operators on
	a second finite-dimensional space. Then
	\[
		\Bigl(\sum_o S_o \ot A_o\Bigr) \Bigl(\sum_o S_o \ot B_o\Bigr)
		= \sum_o S_o \ot (A_o B_o)\;.
	\]
	If in addition $A_o A_o^\dagger = \Id$ for every $o$, then
	$\bigl(\sum_o S_o \ot A_o\bigr) \bigl(\sum_o S_o \ot A_o\bigr)^\dagger
	= \Id$.
\end{lemma}
\begin{proof}
	\leanok
	\uses{def:povm-conventions}
	Expanding the product gives the terms $S_o S_{o'} \ot (A_o B_{o'})$, and
	$S_o S_{o'} = 0$ for $o \ne o'$, so only the diagonal terms
	$S_o \ot (A_o B_o)$ survive. For the second assertion, each $S_o$ is
	self-adjoint, so $\bigl(\sum_o S_o \ot A_o\bigr)^\dagger
	= \sum_o S_o \ot A_o^\dagger$; the first identity then turns the product
	into $\sum_o S_o \ot (A_o A_o^\dagger) = \sum_o S_o \ot \Id = \Id$.
\end{proof}

\begin{definition}[The swap unitaries]\label{def:v-swap-unitary}
	\lean{MIPStarRE.QPBT.swapUnitary}
	\leanok
	\uses{lem:qld-4-7, def:generalized-pauli, def:qld-full-field-decoder}
	Define the linear map
	\[
		V_{\mathrm{A}} = \sum_{g_X,g_Z} (\hat{S}_{g_X,g_Z})_{\mathrm{A}\mathrm{A}'} \ot \bigl(\tau^X(\mathrm{Dec}(g_Z)) \, \tau^Z(\mathrm{Dec}(g_X))\bigr)_{\mathrm{A}''}\;,
	\]
	where $\tau^X(\cdot)$ and $\tau^Z(\cdot)$ are the generalized Pauli observables acting on $(\C^q)^{\ot M}$; note that the argument of $\tau^X$ is $\mathrm{Dec}(g_Z)$ and the argument of $\tau^Z$ is $\mathrm{Dec}(g_X)$. The map $V_{\mathrm{B}}$, acting on registers $\mathrm{B}\mathrm{B}'\mathrm{B}''$, is defined analogously.
\end{definition}

% In the source the following properties are established inline with the
% definition of the swap unitaries; they are recorded here as a lemma.
\begin{lemma}[Unitarity and conjugation action of the swap unitaries]\label{lem:v-swap-conjugation}
	\lean{MIPStarRE.QPBT.swapUnitary_mul_conjTranspose}
	\lean{MIPStarRE.QPBT.conjTranspose_mul_swapUnitary}
	\lean{MIPStarRE.QPBT.swapUnitary_conj_tildeObs}
	\lean{MIPStarRE.QPBT.swapUnitary_conj_tildeM}
	\leanok
	\uses{def:v-swap-unitary, def:tilde-w-observables, def:tilde-m-measurement, def:generalized-pauli, def:indicator-vector, def:bracket}
	The maps $V_{\mathrm{A}}$ and $V_{\mathrm{B}}$ of Definition~\ref{def:v-swap-unitary} are unitaries acting on registers $\mathrm{A}\mathrm{A}'\mathrm{A}''$ and $\mathrm{B}\mathrm{B}'\mathrm{B}''$ respectively, and their conjugation action is exact, as follows. For all $W \in \{X,Z\}$, $j \in \{0,\ldots,t-1\}$, and $\tilde{u} \in \F_q^M$,
	\begin{equation}\label{eq:v-swap-obs-conjugation}
		V_{\mathrm{A}} \, \tilde{W}^j(\tilde{u}) \, V_{\mathrm{A}}^\dagger = \Id_{\mathrm{A}\mathrm{A}'} \ot \bigl( \tau^W(e_j \tilde{u}) \bigr)_{\mathrm{A}''}
		\qquad \text{and} \qquad
		V_{\mathrm{B}} \, \tilde{W}^j(\tilde{u}) \, V_{\mathrm{B}}^\dagger = \Id_{\mathrm{B}\mathrm{B}'} \ot \bigl( \tau^W(e_j \tilde{u}) \bigr)_{\mathrm{B}''}\;,
	\end{equation}
	and for all $W \in \{X,Z\}$, $u \in \F_q^m$, and $a \in \F_q$,
	\begin{equation}\label{eq:qld-unitary-6}
		V_{\mathrm{B}} \, \tilde{M}^{W,\mathrm{ind}_m(u)}_a \, V_{\mathrm{B}}^\dagger = \Id_{\mathrm{B}\mathrm{B}'} \ot (\tau^W_{[g_h(u)=a]})_{\mathrm{B}''}\;,
		\qquad
		\tau^W_{[g_h(u)=a]} = \sum_{h \,:\, g_h(u) = a} \tau^W_h\;,
	\end{equation}
	together with the analogous identity for $V_{\mathrm{A}}$ on registers $\mathrm{A}\mathrm{A}'\mathrm{A}''$.
\end{lemma}
\begin{proof}
	\leanok
	\uses{lem:qld-4-7, lem:tildew-product-form, lem:twisted-commutation, lem:pauli-observable-expansion, lem:controlled-sum-support, lem:distance-tensor-sum-left, def:v-swap-unitary, def:tilde-m-measurement, def:s-w-marginals, def:tau-dot-product-projector, def:indicator-vector, def:qld-full-field-decoder, def:generalized-pauli, def:bracket}
	\leanok
	Since the measurement $\{\hat{S}_{g_X,g_Z}\}$ is projective (Lemma~\ref{lem:qld-4-7}) and the $\tau^X, \tau^Z$ observables are self-inverse,
	\[
		V_{\mathrm{A}} V_{\mathrm{A}}^\dagger = \sum_{g_X,g_Z} (\hat{S}_{g_X,g_Z})_{\mathrm{A}\mathrm{A}'} \ot \Id_{\mathrm{A}''} = \Id_{\mathrm{A}\mathrm{A}'\mathrm{A}''}\;,
	\]
	and the same computation gives $V_{\mathrm{A}}^\dagger V_{\mathrm{A}} = \Id$, so $V_{\mathrm{A}}$ is unitary; likewise $V_{\mathrm{B}}$. For \eqref{eq:v-swap-obs-conjugation}, conjugate the product form \eqref{eq:tildewj-product-form} of $\tilde{W}^j(\tilde{u})$ (Lemma~\ref{lem:tildew-product-form}): by the orthogonality of the projectors $\hat{S}_{g_X,g_Z}$, only the diagonal terms survive, and the twisted commutation relation (Lemma~\ref{lem:twisted-commutation})
	\[
		\tau^X(\mathrm{Dec}(g_Z)) \, \tau^Z(\mathrm{Dec}(g_X)) \, \tau^W(e_j \tilde{u}) = (-1)^{\tr( e_j(\mathrm{Dec}(g_W) \cdot \tilde{u}))} \, \tau^W(e_j \tilde{u}) \, \tau^X(\mathrm{Dec}(g_Z)) \, \tau^Z(\mathrm{Dec}(g_X))
	\]
	cancels the sign $(-1)^{\tr(e_j (\mathrm{Dec}(g_W) \cdot \tilde{u}))}$ term by term, leaving $\sum_{g_X,g_Z} \hat{S}_{g_X,g_Z} \ot \tau^W(e_j \tilde{u}) = \Id \ot \tau^W(e_j \tilde{u})$. For \eqref{eq:qld-unitary-6}, expand $\tilde{M}^{W,\mathrm{ind}_m(u)}_a$ by \eqref{eq:tilde_M} and \eqref{eq:def-tauwu}; by the orthogonality of the $\hat{S}_{g_X,g_Z}$, conjugation replaces, in the term of $\hat{S}^W_g$, each projector $\tau^W_h$ by $\tau^W_{h + \mathrm{Dec}(g)}$, using the identities $\tau^X(s)\, \tau^Z_h \,\tau^X(s)^\dagger = \tau^Z_{h+s}$ and $\tau^Z(s)\, \tau^X_h\, \tau^Z(s)^\dagger = \tau^X_{h+s}$. Substituting $h' = h + \mathrm{Dec}(g)$ turns the constraint $h \cdot \mathrm{ind}_m(u) = (\mathrm{Dec}(g) \cdot \mathrm{ind}_m(u)) - a$ into $h' \cdot \mathrm{ind}_m(u) = -a = a$, since $\F_q$ has characteristic $2$; the quantity $\mathrm{Dec}(g) \cdot \mathrm{ind}_m(u)$ cancels exactly, so no relation between it and $g(u)$ is used (Remark~\ref{rem:qld-decoding-identity}). Since $h' \cdot \mathrm{ind}_m(u) = g_{h'}(u)$ \eqref{eq:low-degree-encoding-definition}, summing over $g$ gives \eqref{eq:qld-unitary-6}; the computation for $V_{\mathrm{A}}$ on registers $\mathrm{A}\mathrm{A}'\mathrm{A}''$ is identical.
\end{proof}

\begin{lemma}[Extraction of the EPR state and the Pauli observables]\label{lem:qld-unitary}
	\lean{MIPStarRE.QPBT.ExtractionWitness}
	\lean{MIPStarRE.QPBT.deltaExtract}
	\lean{MIPStarRE.QPBT.exists_extractionWitness}
	\lean{MIPStarRE.QPBT.exists_extractionWitness_fourth_root_rate}
	\leanok
	\uses{lem:qld-4-7, def:v-swap-unitary, def:expanded-state, lem:projective-strategy-setup, def:generalized-pauli, def:EPR, def:povm-distance}
	There are universal constants $a>1$, $0<b<1$ as in
	Lemma~\ref{lem:qld-4-7}, and $C_{\mathrm{ext}}\geq1$, all chosen before
	the parameters and strategy. Write $\delta_G=a(md)^a
	(\eps^b+q^{-b}+2^{-bmd})$ and set
	\[
		\delta_P^{\mathrm{ext}}=C_{\mathrm{ext}}B,
		\qquad
		\delta_{\mathrm{qld}}=C_{\mathrm{ext}}
		\left((\delta_P^{\mathrm{ext}})^{1/4}+\frac{md}{q}\right).
	\]
	There exist unitaries $V_{\mathrm{A}}$ acting on registers $\mathrm{A}\mathrm{A}'\mathrm{A}''$ and $V_{\mathrm{B}}$ acting on registers $\mathrm{B}\mathrm{B}'\mathrm{B}''$, such that the following hold.
	\begin{enumerate}
		\item There exists a state $\ket{\mathrm{aux}}$ on registers $\mathrm{A}\mathrm{A}'\mathrm{B}\mathrm{B}'$ such that
		\[
			\bigl\| V_{\mathrm{A}} \ot V_{\mathrm{B}} \ket{\hat{\psi}}_{\mathrm{A}\mathrm{A}'\mathrm{A}''\mathrm{B}\mathrm{B}'\mathrm{B}''} - \ket{\mathrm{aux}}_{\mathrm{A}\mathrm{A}'\mathrm{B}\mathrm{B}'} \ot \ket{\mathrm{EPR}_q}^{\ot M}_{\mathrm{A}''\mathrm{B}''} \bigr\|^2 \leq \delta_{\mathrm{qld}}\;.
		\]
		\item For $W \in \{X,Z\}$,
		\begin{align*}
			V_{\mathrm{A}} \, \bigl(M^{(\mathrm{Pauli},W)}_h \bigr)_{\mathrm{A}} \, V_{\mathrm{A}}^\dagger &\approx_{\delta_{\mathrm{qld}}} \Id_{\mathrm{A}\mathrm{A}'} \ot (\tau^W_h)_{\mathrm{A}''}\;, \\
			V_{\mathrm{B}} \, \bigl(M^{(\mathrm{Pauli},W)}_h \bigr)_{\mathrm{B}} \, V_{\mathrm{B}}^\dagger &\approx_{\delta_{\mathrm{qld}}} \Id_{\mathrm{B}\mathrm{B}'} \ot (\tau^W_h)_{\mathrm{B}''}\;,
		\end{align*}
		where the $\approx$ statements hold with respect to the state $\ket{\mathrm{aux}}_{\mathrm{A}\mathrm{A}'\mathrm{B}\mathrm{B}'} \ot \ket{\mathrm{EPR}_q}^{\ot M}_{\mathrm{A}''\mathrm{B}''}$, with the answer summation over $h \in \F_q^{M}$ and no average over questions.
	\end{enumerate}
	The unitaries may be taken to be $V_{\mathrm{A}}$ and $V_{\mathrm{B}}$ of Definition~\ref{def:v-swap-unitary}.
	For every admissible $(q,m,d)$, every $0\leq\eps\leq1$, and every projective
	strategy with success probability at least $1-\eps$, the global measurement
	of Lemma~\ref{lem:qld-4-7} and the above extraction data exist together.
	The conditional result below proves the extraction estimates once that
	measurement is fixed; it is not an additional hypothesis here. The subsequent
	range-projection transfer to isometries remains separate, as recorded in
	\cite{gap:qpbt_extraction-transfer}.
	Moreover, the universal constants can be chosen so that the same global
	measurement has consistency error at most $\delta_G$ and both extraction
	conclusions hold at the source's fourth-root scale
	$C_{\mathrm{ext}}(\delta_G^{1/4}+md/q)$, with the \emph{same} $\delta_G$
	in both bounds. This stronger numerical assertion is separate from the
	nested construction/extraction estimate displayed above. It retains the
	directly indexed low-degree carrier of Lemma~\ref{lem:qld-4-7}; the
	seed-indexed correspondence is not established by this assertion.
\end{lemma}
\begin{proof}
	\leanok
	\uses{lem:qld-4-7, lem:qld-unitary-given-global-pair}
	Choose the universal constants in Lemma~\ref{lem:qld-4-7} and the
	conditional extraction estimate below before fixing the parameters and
	strategy. The former supplies both projective polynomial-pair measurements
	at error $\delta_G$. Since $\delta_G\geq0$, the latter applies to this same
	pair and yields the concrete swap unitaries and the normalized auxiliary state
	at the displayed scale. In particular, the global measurement is an output
	of the argument, not an assumed input.
	For the fourth-root form, first bound the nested error by a second member
	of the $\delta_G$ family, then increase its leading constant and decrease
	its exponent so it also bounds the original global-measurement error.
	Weaken only the two consistency certificates, leaving the measurements
	and swap maps unchanged. When the enlarged error $x$ is at most one,
	$\delta_{\mathrm{ext}}\leq x\leq x^{1/4}$; when $x>1$, completeness and
	unitarity give both extraction estimates at error four, which is at most
	$C_{\mathrm{ext}}(x^{1/4}+md/q)$ after enlarging $C_{\mathrm{ext}}$.
\end{proof}

\begin{theorem}[EPR projection estimates]\label{thm:qld-epr-projection-estimates}
	\lean{MIPStarRE.QPBT.Extraction.pauliCorrelation_mul_eq_epr}
	\lean{MIPStarRE.QPBT.Extraction.norm_sub_eprProjection_le}
	\lean{MIPStarRE.QPBT.ProjectiveSetting.idealExpState_extractionAux0}
	\lean{MIPStarRE.QPBT.ProjectiveSetting.exists_unit_aux_near_eprProjection}
	\leanok
	\uses{def:generalized-pauli, def:EPR, lem:projective-strategy-setup}
	This auxiliary statement isolates the projection and normalization
	calculations in the proof of Lemma~\ref{lem:qld-unitary}.
	Write $E=\ket{\mathrm{EPR}_q}^{\ot M}$ and
	\[
		H_W=\frac{1}{q^M}\sum_{u\in\F_q^M}\tau^W(u)\ot\tau^W(u).
	\]
	Then $H_XH_Z=\ket{E}\bra{E}$. Let $R$ be an auxiliary register and
	$\vartheta$ a unit vector on $R\ot A''\ot B''$. If $\delta\geq0$ and
	\[
		\Re\langle\vartheta,(\Id_R\ot H_W)\vartheta\rangle
		\geq1-\delta/2\qquad(W\in\{X,Z\}),
	\]
	then, writing $\Pi=\Id_R\ot\ket{E}\bra{E}$, one has
	\[
		\|\vartheta-\Pi\vartheta\|\leq2\sqrt\delta.
	\]
	In the projective setting of Lemma~\ref{lem:projective-strategy-setup},
	take $R=AA'BB'$. For every vector $\vartheta$ on the six registers, set
	$\mathrm{aux}_0=(\Id_R\ot\bra{E})\vartheta$. Then
	$\mathrm{aux}_0\ot E=\Pi\vartheta$. If $\vartheta$ is a unit vector, there is
	a unit vector $\mathrm{aux}$ on $R$ such that
	\[
		\|\vartheta-\mathrm{aux}\ot E\|
		\leq2\|\vartheta-\Pi\vartheta\|.
	\]
\end{theorem}
\begin{proof}
	\leanok
	\uses{lem:pauli-observable-expansion, def:generalized-pauli, def:EPR}
	Fourier inversion shows that $H_Z$ keeps precisely the equal computational
	labels. Averaging the simultaneous shifts gives $H_XH_Z=\ket{E}\bra{E}$.
	Each $\Id_R\ot H_W$ is an average of unitaries, hence is a contraction.
	Expanding the squared distance gives
	$\|\vartheta-(\Id_R\ot H_W)\vartheta\|^2\leq\delta$.
	The product identity and the triangle inequality now bound
	$\|\vartheta-\Pi\vartheta\|$ by $2\sqrt\delta$.
	The partial-inner-product formula gives
	$\mathrm{aux}_0\ot E=\Pi\vartheta$. Normalizing $\mathrm{aux}_0$
	at most doubles its distance from $\vartheta$, by the reverse triangle
	inequality. If $\mathrm{aux}_0=0$, the projection error is one; the
	original strategy state tensored with one EPR pair supplies a unit
	reference vector on $R$, and the required distance is at most two.
\end{proof}

\begin{example}[State extraction from a supplied global witness]
	\label{lem:qld-state-extraction-given-global-pair}
	\lean{MIPStarRE.QPBT.exists_extraction_aux_ofGlobalPairWitness}
	\leanok
	\uses{lem:qld-construct-the-paulis-given-global-pair, def:v-swap-unitary,
		thm:qld-epr-projection-estimates}
	Assume the global polynomial-pair measurement and both consistency estimates
	of Lemma~\ref{lem:qld-4-7} as supplied data at error $\delta_G\geq0$.
	There is a universal constant $C\geq1$ such that, for every admissible
	projective setting with $0\leq\eps\leq1$, a normalized auxiliary state satisfies
	\[
	  \|V_{\mathrm A}\ot V_{\mathrm B}\ket{\hat\psi}
	      -\ket{\mathrm{aux}}\ot\ket{\mathrm{EPR}_q}^{\ot M}\|^2
	  \leq16C(\delta_G+\sqrt\eps+md/q).
	\]
	This statement assumes the global measurement and therefore does not assert
	its construction or the unconditional extraction conclusion.
\end{example}

\begin{example}[Evaluated Pauli consistency from a supplied global witness]
  \label{lem:qld-evaluated-pauli-given-global-pair}
  \lean{MIPStarRE.QPBT.tildeM_consistent_pointMeas_ofGlobalPairWitness}
  \lean{MIPStarRE.QPBT.tildeM_consistent_pointMeas'_ofGlobalPairWitness}
  \lean{MIPStarRE.QPBT.evaluated_pauli_tilde_consistency_ofGlobalPairWitness}
  \leanok
  \uses{lem:qld-construct-the-paulis-given-global-pair, lem:qld-win-implications}
  Assume the global polynomial-pair witness at error $\delta_G\geq0$ as data.
  Each of the two point comparisons in Item 1 of
  Lemma~\ref{lem:qld-construct-the-paulis} has a universal bound $CB$, where
  $B=\delta_G+\sqrt\eps+md/q$ and $0\leq\eps\leq1$.
  After enlarging the universal constant, the two consistency defects between
  the evaluated total-Pauli measurement of one player and the pulled-apart
  measurement of the other player are also at most $CB$ on $\ket{\hat\psi}$.
  The constants precede all parameters, strategies, and supplied witnesses.
  No construction of the global witness is asserted.
\end{example}

\begin{lemma}[Pauli comparison after a change of state]
  \label{lem:qld-pauli-overlap-state-transfer}
  \lean{MIPStarRE.QPBT.pauli_overlap_transfer}
  \lean{MIPStarRE.QPBT.pauli_distance_on_ideal_le}
  \lean{MIPStarRE.QPBT.extraction_small_error_absorption}
  \leanok
  \uses{def:generalized-pauli, def:EPR, def:povm-distance}
  Let $A=\{A_h\}$ and $B=\{B_h\}$ be arbitrary local POVMs, and let
  $\vartheta$ and $\Phi$ be unit vectors on the associated bipartite space.
  Write
  \[
    E_{A,B}(\vartheta)=\mathbb E_u\sum_a
      \langle\vartheta,A_{[g_h(u)=a]}\ot B_{[g_h(u)=a]}\vartheta\rangle.
  \]
  Then
  \[
    E_{A,B}(\vartheta)
    \leq \sum_h\langle\Phi,A_h\ot B_h\Phi\rangle
      +md/q+2\|\vartheta-\Phi\|.
  \]
  In particular, let $A$ act on
  $\mathcal H_A\ot\mathbb C^{q^M}$, take $B$ to be the canonical Pauli
  projectors on $\mathcal H_B\ot\mathbb C^{q^M}$, and set
  $\Delta=\mathrm{aux}\ot\mathrm{EPR}_q^{\ot M}$ in local player order for a
  unit auxiliary vector. For either Pauli basis,
  \[
    \sum_h\|((A_h-\Id_{\mathcal H_A}\ot\tau^W_h)\ot\Id)\Delta\|^2
    \leq 2(1-E_{A,B}(\vartheta))+2md/q+4\|\vartheta-\Delta\|.
  \]
  These bounds are independent of the dimensions and the number of answers.
  Finally, for $0\leq\delta\leq1$ and $C\geq18$, both $16\delta$ and
  $2\delta+2md/q+16\sqrt\delta$ are at most $C(\delta^{1/4}+md/q)$.
\end{lemma}
\begin{proof}
  \leanok
  \uses{lem:sandwich-evaluated-diagonal-overlap,
    lem:schwartz-zippel-total-degree, fact:agreement,
    lem:pauli-observable-expansion}
  A selected sum of effects of the product POVM is a positive contraction.
  Its expectation changes by at most twice the distance between unit vectors.
  Schwartz--Zippel bounds the unequal-encoding contribution by $md/q$.
  Symmetry of the Pauli projectors moves their action across the EPR pair,
  and the agreement inequality bounds the full distance by twice the overlap
  deficit. The scalar bounds follow from $\delta,\sqrt\delta\leq\delta^{1/4}$.
\end{proof}

\begin{example}[Extraction with allowed error at least four]
	\label{lem:qld-large-error-extraction-given-global-pair}
	\lean{MIPStarRE.QPBT.exists_extractionWitness_ofGlobalPairWitness_of_four_le}
	\leanok
	\uses{def:v-swap-unitary, def:EPR, lem:projective-strategy-setup}
	For any supplied global polynomial-pair witness and any allowed squared
	error $\delta\geq4$, there exists a conditional extraction witness.
	The auxiliary vector is the original strategy state tensored with EPR on
	the auxiliary registers. Unitarity bounds the state distance by four;
	completeness bounds the summed Pauli measurement distance by four, for
	each player and each basis. No restriction on a projection norm is needed.
\end{example}

\begin{example}[Conditional extraction given a global polynomial-pair witness]
	\label{lem:qld-unitary-given-global-pair}
	\lean{MIPStarRE.QPBT.exists_extractionWitness_ofGlobalPairWitness}
	\leanok
	\uses{lem:qld-4-7, def:v-swap-unitary, def:expanded-state,
		lem:projective-strategy-setup, def:generalized-pauli, def:EPR,
		def:povm-distance}
	This formalization-only auxiliary is not a named statement of the paper. It is
	the conditional, given-witness form
	of Lemma~\ref{lem:qld-unitary}: the projective global polynomial-pair measurement
	$\{\hat{S}_{g_X,g_Z}\}$ of Lemma~\ref{lem:qld-4-7}, together with the two
	point-consistency estimates that lemma asserts at a parameter $\delta_G \geq 0$,
	is assumed as data rather than obtained from that lemma. There is a universal
	constant $C_{\mathrm{ext}} \geq 1$ such that the following holds for every
	admissible $(q,m,d)$, every $0 \leq \eps \leq 1$, every projective strategy
	setup at error $\eps$ (Lemma~\ref{lem:projective-strategy-setup}), every
	$\delta_G \geq 0$, and every assumed measurement as above. The swap unitaries
	$V_{\mathrm{A}}$ and $V_{\mathrm{B}}$ built from that measurement by
	Definition~\ref{def:v-swap-unitary} are two-sided unitaries, and there is a
	normalized state $\ket{\mathrm{aux}}$ on registers
	$\mathrm{A}\mathrm{A}'\mathrm{B}\mathrm{B}'$ for which items 1 and 2 of
	Lemma~\ref{lem:qld-unitary} hold at the scale $\delta_{\mathrm{qld}}$ of that
	lemma, computed from $\eps$ and the assumed $\delta_G$.
	This statement assumes the global measurement and therefore does not assert
	its construction or the unconditional conclusion of
	Lemma~\ref{lem:qld-unitary}. The source-facing composition uses this
	conditional result after constructing the measurement; see
	\cite{gap:qpbt_extraction-transfer}.
\end{example}
\begin{proof}
	\uses{lem:qld-state-extraction-given-global-pair,
		lem:qld-evaluated-pauli-given-global-pair, lem:v-swap-conjugation,
		lem:qld-pauli-overlap-state-transfer,
		lem:qld-large-error-extraction-given-global-pair}
	Let $C_S,C_E$ be the universal constants in the auxiliary-state and
	evaluated-consistency estimates. Choose
	$C_{\mathrm{ext}}=\max\{C_S,C_E,18\}$ before fixing any parameters or
	strategy, and put $\delta=C_{\mathrm{ext}}(\delta_G+\sqrt\eps+md/q)$.
	The normalized auxiliary state has squared extraction error at most
	$16\delta$. For each player and either Pauli basis, simultaneous swap
	conjugation and \eqref{eq:qld-unitary-6} identify the evaluated overlap
	on the swapped state with one minus the corresponding consistency defect
	on $\ket{\hat\psi}$. That defect is at most $\delta$.
	Lemma~\ref{lem:qld-pauli-overlap-state-transfer} therefore bounds the
	full Pauli distance on the ideal state by
	$2\delta+2md/q+16\sqrt\delta$.
	When $\delta\leq1$, both errors are at most
	$C_{\mathrm{ext}}(\delta^{1/4}+md/q)$ by the scalar estimate in that lemma.
	When $\delta>1$, the allowed error is at least four, so the complete
	large-error construction applies. The exact swap identities give both
	unitarity equations in either case.
\end{proof}

\begin{lemma}[Preservation of the robustness error family]\label{lem:qld-extraction-error-form}
	\lean{MIPStarRE.QPBT.deltaExtract_le_deltaQld}
	\leanok
	\uses{def:admissible, lem:qld-unitary}
	Let $C\geq1$, $a>1$, and $0<b<1$. Suppose that
	\[
		\delta_G=a(md)^a\bigl(\eps^b+q^{-b}+2^{-bmd}\bigr),
		\qquad
		\delta_P=C\left(\delta_G+\sqrt{\eps}+\frac{md}{q}\right),
	\]
	and set $\delta_{\mathrm{ext}}=C(\delta_P^{1/4}+md/q)$. There are constants
	$a'\geq1$ and $0<b'<1$, depending only on $C,a,b$, such that, for every
	admissible $(q,m,d)$ and every $0\leq\eps\leq1$,
	\[
		\delta_{\mathrm{ext}}
		\leq a'(md)^{a'}\bigl(\eps^{b'}+q^{-b'}+2^{-b'md}\bigr).
	\]
\end{lemma}
\begin{proof}
	\leanok
	\uses{def:admissible}
	Write $D=md$, and set $\beta=b/8$, $A=4C^2a$, $H=aD^a$, and
	$R=\eps^\beta+q^{-\beta}+2^{-\beta D}$.
	Admissibility gives $D,q\geq1$, so $H\geq D\geq1$, $A\geq a>1$,
	and $0<\beta<1$. Fourth-root subadditivity and monotonicity of real
	powers, using $0\leq\eps\leq1$ and $\beta\leq b/4$, give
	\[
		\delta_G^{1/4}\leq HR,\qquad
		\eps^{1/8}\leq HR,\qquad
		(D/q)^{1/4}\leq HR,\qquad D/q\leq HR.
	\]
	Indeed, the fourth root of the sum defining $\delta_G/H$ is at most
	$R$, while $H^{1/4}\leq H$, $D^{1/4}\leq D\leq H$,
	$\eps^{1/8}\leq\eps^\beta$, and $q^{-1/4},q^{-1}\leq q^{-\beta}$.
	Since $C^{1/4}\leq C$, subadditivity also gives
	$\delta_P^{1/4}\leq3CHR$. Consequently,
	\[
		\delta_{\mathrm{ext}}
		\leq C(3C+1)HR\leq4C^2HR=AD^aR\leq AD^AR.
	\]
	Thus $a'=A$ and $b'=\beta$ have the required properties.
\end{proof}

\begin{remark}[Two steps in the source proof of the extraction lemma]\label{rem:qld-unitary-triangle-slip}
	In the source proof of Lemma~\ref{lem:qld-unitary}, the triangle inequality gives $\| H_X \ket{\vartheta} - H_Z \ket{\vartheta} \| \leq \| (\Id - H_X)\ket{\vartheta} \| + \| (\Id - H_Z)\ket{\vartheta} \| \leq 2\sqrt{\delta_P}$, so that $\| H_X \ket{\vartheta} - H_Z \ket{\vartheta} \|^2 \leq 4\delta_P$; the source instead asserts the bound $2\sqrt{\delta_P}$ for the \emph{squared} norm, which holds only when $\sqrt{\delta_P} \leq 1/2$, and consequently carries the term $\sqrt{\delta_P}$ into an overlap bound of $1 - 2\delta_P - \sqrt{\delta_P}$. The proof above uses the corrected bound $4\delta_P$, which gives the overlap bound $1 - 4\delta_P$ and leaves the final bound $\delta_{\mathrm{qld}} = O(\delta_P^{1/4} + md/q)$ unchanged. The source proof also normalizes the vector $(\bra{\mathrm{EPR}_q}^{\ot M})_{\mathrm{A}''\mathrm{B}''} \ket{\vartheta}$ without ensuring that it is nonzero; the case distinction on $\delta_P$ in the proof above supplies this, the conclusions of the lemma being trivial when $\delta_P$ exceeds a universal constant.
	The numerical corrections and their role in the extraction argument are
	documented in \cite{gap:qpbt_extraction-transfer}.
\end{remark}

\section{Proof of the soundness of the Pauli basis test}

\subsection{Range projections and isometry conjugation}

The following formalization-only auxiliary results are not named theorems of
the source article. They justify the insertion of range projections in the
final soundness argument. The omitted step in the source is discussed in
\cite{gap:qpbt_extraction-transfer}.

\begin{theorem}[Range-projection estimate]\label{thm:pauli-range-projection-support}
  \lean{MIPStarRE.QPBT.sum_norm_mul_projection_sub_sq_le}
  \leanok
  Let $\{N_a\}_{a\in A}$ be a complete POVM on a finite-dimensional complex
  Hilbert space $\mathcal H$, with $A$ finite, and let $\{T_a\}_{a\in A}$ be
  any operator family on $\mathcal H$. Let $Q$ be an orthogonal projection
  and let $\theta,\Delta\in\mathcal H$ satisfy $Q\theta=\theta$. Then
  \[
    \sum_a\|(N_aQ-T_a)\Delta\|^2
    \le 2\sum_a\|(N_a-T_a)\Delta\|^2+2\|\Delta-\theta\|^2.
  \]
\end{theorem}
\begin{proof}
  \uses{lem:sandwich-povm-square-sum}
  \leanok
  Since $I-Q$ is a contraction and annihilates $\theta$,
  $\|(I-Q)\Delta\|\le\|\Delta-\theta\|$.
  Decompose $(N_aQ-T_a)\Delta=(N_a-T_a)\Delta-N_a(I-Q)\Delta$
  and use $\|x-y\|^2\le2\|x\|^2+2\|y\|^2$, a consequence of the
  parallelogram identity. Summing the second terms uses
  $\sum_a N_a^\dagger N_a\le I$ and gives the asserted estimate.
\end{proof}

\begin{theorem}[Isometry transfer]\label{thm:pauli-isometry-transfer-support}
  \lean{MIPStarRE.QPBT.sum_norm_conjIsometry_sub_sq_le}
  \leanok
  Let $\phi\colon\mathcal H\to\mathcal K$ be a linear isometry between
  finite-dimensional complex Hilbert spaces. Let $\{M_a\}_{a\in A}$ be
  operators on $\mathcal H$, let $\{N_a\}_{a\in A}$ be a complete POVM on
  $\mathcal K$, and let $\{T_a\}_{a\in A}$ be arbitrary operators on
  $\mathcal K$, with $A$ finite. Suppose $N_a\phi=\phi M_a$ for every $a$.
  For all $\psi\in\mathcal H$ and $\Delta\in\mathcal K$,
  \[
    \sum_a\|(\phi M_a\phi^\dagger-T_a)\Delta\|^2
    \le2\sum_a\|(N_a-T_a)\Delta\|^2+2\|\Delta-\phi\psi\|^2.
  \]
\end{theorem}
\begin{proof}
  \uses{thm:pauli-range-projection-support}
  \leanok
  Set $Q=\phi\phi^\dagger$. It is an orthogonal projection fixing
  $\phi\psi$, and the intertwining identity gives
  $N_aQ=\phi M_a\phi^\dagger$. Apply
  Theorem~\ref{thm:pauli-range-projection-support}.
\end{proof}

\begin{theorem}[Alice's bipartite isometry transfer]
  \label{thm:pauli-isometry-transfer-alice-support}
  \lean{MIPStarRE.QPBT.sum_norm_leftTensor_conjIsometry_sub_sq_le}
  \leanok
  Let $\phi_{\mathrm A}\colon\mathcal H_{\mathrm A}\to\mathcal K_{\mathrm A}$
  and $\phi_{\mathrm B}\colon\mathcal H_{\mathrm B}\to\mathcal K_{\mathrm B}$
  be linear isometries between finite-dimensional complex Hilbert spaces.
  Let $A$ be finite, let $\{M_a\}_{a\in A}$ be operators on
  $\mathcal H_{\mathrm A}$, and let $\{N_a\}_{a\in A}$ be a complete POVM
  on $\mathcal K_{\mathrm A}$ satisfying $N_a\phi_{\mathrm A}=\phi_{\mathrm A}M_a$.
  For any operators $T_a$ on $\mathcal K_{\mathrm A}\ot\mathcal K_{\mathrm B}$
  and vectors $\psi\in\mathcal H_{\mathrm A}\ot\mathcal H_{\mathrm B}$,
  $\Delta\in\mathcal K_{\mathrm A}\ot\mathcal K_{\mathrm B}$, writing
  $\theta=(\phi_{\mathrm A}\ot\phi_{\mathrm B})\psi$, we have
  \[
    \sum_a\|((\phi_{\mathrm A}M_a\phi_{\mathrm A}^\dagger)\ot I-T_a)\Delta\|^2
    \le2\sum_a\|(N_a\ot I-T_a)\Delta\|^2+2\|\Delta-\theta\|^2.
  \]
\end{theorem}
\begin{proof}
  \uses{thm:pauli-range-projection-support}
  \leanok
  Apply Theorem~\ref{thm:pauli-range-projection-support} to the complete POVM
  $\{N_a\ot I\}$ and projection $Q=(\phi_{\mathrm A}\phi_{\mathrm A}^\dagger)\ot I$.
  The projection fixes $\theta$, and
  $(N_a\ot I)Q=(\phi_{\mathrm A}M_a\phi_{\mathrm A}^\dagger)\ot I$.
\end{proof}

\begin{theorem}[Bob's bipartite isometry transfer]
  \label{thm:pauli-isometry-transfer-bob-support}
  \lean{MIPStarRE.QPBT.sum_norm_rightTensor_conjIsometry_sub_sq_le}
  \leanok
  Let $\phi_{\mathrm A}\colon\mathcal H_{\mathrm A}\to\mathcal K_{\mathrm A}$
  and $\phi_{\mathrm B}\colon\mathcal H_{\mathrm B}\to\mathcal K_{\mathrm B}$
  be linear isometries between finite-dimensional complex Hilbert spaces.
  Let $A$ be finite, let $\{M_a\}_{a\in A}$ be operators on
  $\mathcal H_{\mathrm B}$, and let $\{N_a\}_{a\in A}$ be a complete POVM
  on $\mathcal K_{\mathrm B}$ satisfying
  $N_a\phi_{\mathrm B}=\phi_{\mathrm B}M_a$ for every $a$.
  For arbitrary operators $T_a$ on
  $\mathcal K_{\mathrm A}\ot\mathcal K_{\mathrm B}$ and vectors
  $\psi\in\mathcal H_{\mathrm A}\ot\mathcal H_{\mathrm B}$,
  $\Delta\in\mathcal K_{\mathrm A}\ot\mathcal K_{\mathrm B}$,
  write $\theta=(\phi_{\mathrm A}\ot\phi_{\mathrm B})\psi$. Then
  \[
    \sum_a\|(I\ot(\phi_{\mathrm B}M_a\phi_{\mathrm B}^\dagger)-T_a)\Delta\|^2
    \le2\sum_a\|(I\ot N_a-T_a)\Delta\|^2+2\|\Delta-\theta\|^2.
  \]
\end{theorem}
\begin{proof}
  \uses{thm:pauli-range-projection-support}
  \leanok
  Apply Theorem~\ref{thm:pauli-range-projection-support} to $\{I\ot N_a\}$ and
  $Q=I\ot(\phi_{\mathrm B}\phi_{\mathrm B}^\dagger)$, which fixes $\theta$.
  Intertwining gives $(I\ot N_a)Q=I\ot(\phi_{\mathrm B}M_a\phi_{\mathrm B}^\dagger)$.
\end{proof}

\begin{theorem}[Ancilla isometry and intertwining]
  \label{thm:pauli-ancilla-intertwining-support}
  \lean{MIPStarRE.QPBT.ancillaBlockIsometry}
  \lean{MIPStarRE.QPBT.ancillaBlockIsometry_matrix_apply}
  \lean{MIPStarRE.QPBT.ancillaBlockIsometry_intertwines}
  \lean{MIPStarRE.QPBT.unitaryAncillaIsometry}
  \lean{MIPStarRE.QPBT.unitaryAncillaIsometry_matrix}
  \lean{MIPStarRE.QPBT.unitaryAncillaIsometry_intertwines}
  \leanok
  Let $\eta\in\mathcal K\ot\mathcal L$ be a unit vector, and define
  $J\colon\mathcal H\to(\mathcal H\ot\mathcal K)\ot\mathcal L$ by
  $J\psi=\psi\ot\eta$, with the indicated reassociation of factors.
  Then $J$ is an isometry and, for every operator $M$ on $\mathcal H$,
  $(M\ot I\ot I)J=JM$.
  If $U^\dagger U=I$ on the enlarged space, then $\phi=UJ$ is an isometry
  and
  \[
    \bigl(U(M\ot I\ot I)U^\dagger\bigr)\phi=\phi M.
  \]
  These identities hold for arbitrary finite-dimensional spaces. They give
  the exact intertwining relation used in the final soundness proof when
  $\eta$ is the local EPR state and $U$ is the swap unitary.
\end{theorem}
\begin{proof}
  \leanok
  The tensor norm gives $\|J\psi\|=\|\psi\|\|\eta\|=\|\psi\|$.
  The operator $M$ acts on the first tensor factor, so
  $(M\ot I\ot I)J\psi=J(M\psi)$.
  Finally, $\phi^\dagger\phi=J^\dagger U^\dagger UJ=I$, and
  $U(M\ot I\ot I)U^\dagger UJ=UJM$.
\end{proof}

\begin{theorem}[Formalization support for the extraction-data isometries]
  \label{thm:pauli-extraction-isometry-construction-support}
  \lean{MIPStarRE.QPBT.ExtractionWitness.toPauliSoundnessWitness}
  \lean{MIPStarRE.QPBT.isometryTensor_ancilla_eq_psiHat}
  \lean{MIPStarRE.QPBT.ExtractionWitness.isometryTensor_eq}
  \lean{MIPStarRE.QPBT.ExtractionWitness.idealState_eq}
  \lean{MIPStarRE.QPBT.ExtractionWitness.state_error_eq}
  \leanok
  \uses{lem:qld-unitary, def:expanded-state, def:EPR}
  This node is not a named statement of the paper; it isolates the register
  bookkeeping that the final appendix argument performs silently and that the
  formal proof must carry out explicitly. Let $S$ be a projective Pauli-test
  setting with original local spaces $\mathcal H_{\mathrm A}$ and
  $\mathcal H_{\mathrm B}$, let $w$ be a global polynomial-pair witness for
  $S$, and let $v$ be an extraction witness for $S,w$ with squared error at
  most $\delta$. Adjoining the normalized EPR ancilla to each original local
  space and applying the swap unitary supplied by $v$ gives two local
  isometries, which together with the normalized auxiliary state of $v$ form
  the data compared by the test. For these data the following identities hold,
  where the six registers are regrouped as
  $(\mathrm{A}\mathrm{A}')\mathrm{A}''$ and
  $(\mathrm{B}\mathrm{B}')\mathrm{B}''$.
  \begin{enumerate}
    \item The tensor of the two plain ancilla embeddings applied to the
      strategy state is the expanded state of $S$.
    \item The tensor of the two isometries applied to the strategy state is
      the simultaneous swap action on the expanded state.
    \item The ideal comparison state of the test built from the auxiliary
      state of $v$ is the ideal expanded state of $v$.
    \item Consequently the squared state error compared by the test equals the
      squared error of the supplied extraction data.
  \end{enumerate}
  No success, consistency, or range hypothesis is used; only the identification
  of registers.
\end{theorem}
\begin{proof}
  \leanok
  \uses{thm:pauli-ancilla-intertwining-support, def:expanded-state}
  The first identity is a coefficientwise computation with the tensor
  expansion of the EPR ancilla. The second follows from it, because applying
  the two swap unitaries commutes with the regrouping of registers. The third
  is the same regrouping applied to the auxiliary and EPR factors, and the
  fourth follows from the first three because the regrouping preserves norms.
\end{proof}

\begin{example}[Formalization support: state component of the transfer]
  \label{lem:pauli-extraction-state-distance-support}
  \lean{MIPStarRE.QPBT.ExtractionWitness.state_close_ofExtractionWitness}
  \leanok
  \uses{lem:qld-unitary, def:expanded-state}
  This node is not a named statement of the paper; it records the state
  component of the final appendix estimate for supplied extraction data. In
  the situation of
  Theorem~\ref{thm:pauli-extraction-isometry-construction-support}, the
  squared distance between the transferred strategy state and the ideal state
  is at most $\delta$. Since the left-hand side is the square of a norm, this
  is exactly the bound $\sqrt\delta$ on the state distance itself.
\end{example}
\begin{proof}
  \leanok
  \uses{thm:pauli-extraction-isometry-construction-support}
  Rewrite the squared error through the register identification and apply the
  state estimate carried by the extraction witness.
\end{proof}

\begin{theorem}[Formalization support: Alice's operator component of the transfer]
  \label{thm:pauli-extraction-alice-distance-support}
  \lean{MIPStarRE.QPBT.ExtractionWitness.pauli_distance_alice_le}
  \leanok
  \uses{lem:qld-unitary, def:povm-distance}
  This node is not a named statement of the paper; it records Alice's operator
  component of the final appendix estimate for supplied extraction data. In
  the situation of
  Theorem~\ref{thm:pauli-extraction-isometry-construction-support}, for every
  $W\in\{X,Z\}$ the complete squared Pauli operator-family distance for Alice,
  summed over all $h\in\F_q^M$ and evaluated on the ideal state, is at most
  $4\delta$.
\end{theorem}
\begin{proof}
  \leanok
  \uses{thm:pauli-isometry-transfer-alice-support,
    thm:pauli-ancilla-intertwining-support,
    thm:pauli-extraction-isometry-construction-support,
    lem:pauli-extraction-state-distance-support}
  The swapped total-Pauli family is a complete POVM and satisfies the exact
  intertwining identity for the ancilla isometry, so Alice's bipartite
  isometry transfer applies. Reindexing the extraction witness's own Pauli
  comparison into the test register order bounds the first term by $2\delta$,
  and the state component bounds the second term by $2\delta$.
\end{proof}

\begin{theorem}[Formalization support: Bob's operator component of the transfer]
  \label{thm:pauli-extraction-bob-distance-support}
  \lean{MIPStarRE.QPBT.ExtractionWitness.pauli_distance_bob_le}
  \leanok
  \uses{lem:qld-unitary, def:povm-distance}
  This node is not a named statement of the paper; it records Bob's operator
  component of the same estimate. In the situation of
  Theorem~\ref{thm:pauli-extraction-isometry-construction-support}, for every
  $W\in\{X,Z\}$ the complete squared Pauli operator-family distance for Bob,
  summed over all $h\in\F_q^M$ and evaluated on the ideal state, is at most
  $4\delta$. The two original local spaces are never identified with one
  another.
\end{theorem}
\begin{proof}
  \leanok
  \uses{thm:pauli-isometry-transfer-bob-support,
    thm:pauli-ancilla-intertwining-support,
    thm:pauli-extraction-isometry-construction-support,
    lem:pauli-extraction-state-distance-support}
  The argument is the one for Alice, with Bob's bipartite isometry transfer in
  place of Alice's.
\end{proof}

\begin{theorem}[Formalization support: the transferred state lies in both isometry ranges]
  \label{thm:pauli-extraction-transferred-range-support}
  \lean{MIPStarRE.QPBT.ExtractionWitness.transformed_state_in_ranges}
  \leanok
  \uses{lem:qld-unitary}
  This node is not a named statement of the paper; it records the range
  identities that the appendix argument uses when it says that the transferred
  state lies in the range of each local isometry. In the situation of
  Theorem~\ref{thm:pauli-extraction-isometry-construction-support}, the range
  projection of Alice's isometry, and likewise that of Bob's, fixes the
  transferred strategy state. No range condition is imposed on the ideal
  comparison state.
\end{theorem}
\begin{proof}
  \leanok
  \uses{thm:pauli-extraction-isometry-construction-support}
  Each range projection acts on its own tensor factor as the conjugation of
  the identity by that isometry, and that operator fixes every vector in the
  image of the isometry tensor.
\end{proof}

\begin{theorem}[Concrete Pauli isometry transfer from extraction data]
  \label{thm:pauli-concrete-isometry-transfer-support}
  \lean{MIPStarRE.QPBT.ExtractionWitness.isometry_transfer_bounds}
  \leanok
  \uses{lem:qld-unitary, def:povm-distance}
  Let $S$ be a projective Pauli-test setting with original local spaces
  $\mathcal H_{\mathrm A}$ and $\mathcal H_{\mathrm B}$. Suppose $w$ is a
  global polynomial-pair witness for $S$, and $v$ is an extraction witness
  for $S,w$ with squared error at most $\delta$. Then there exist local
  isometries and a normalized auxiliary state whose state norm distance is at
  most $\max\{\sqrt\delta,4\delta\}$ and whose complete squared Pauli
  operator-family distance on the ideal state is at most
  $\max\{\sqrt\delta,4\delta\}$ for Alice and for Bob, for every
  $W\in\{X,Z\}$. This statement asserts only the common bound and only the
  existence of such data. The sharper component bounds $\sqrt\delta$ and
  $4\delta$, and the explicit isometries obtained from the EPR ancilla and the
  swap unitaries of $v$ that realize them, are recorded separately in the
  subordinate statements~\ref{thm:pauli-extraction-isometry-construction-support},
  \ref{lem:pauli-extraction-state-distance-support},
  \ref{thm:pauli-extraction-alice-distance-support},
  and~\ref{thm:pauli-extraction-bob-distance-support}.
\end{theorem}
\begin{proof}
  \uses{thm:pauli-extraction-isometry-construction-support,
    lem:pauli-extraction-state-distance-support,
    thm:pauli-extraction-alice-distance-support,
    thm:pauli-extraction-bob-distance-support}
  \leanok
  Take the isometries and auxiliary state built from the EPR ancilla and the
  swap unitaries of $v$ as the witness. The state bound follows from the
  squared state component by taking square roots, and the two operator bounds
  are the Alice and Bob components. Each of the three is then relaxed to the
  common maximum.
\end{proof}

\begin{example}[Formalization support: state bound in the robustness error family]
  \label{lem:pauli-extraction-state-error-form-support}
  \lean{MIPStarRE.QPBT.extraction_state_norm_le_deltaQld_ofExtractionWitness}
  \leanok
  \uses{lem:qld-unitary, lem:qld-extraction-error-form}
  This node is not a named statement of the paper; it records the state
  estimate alone in the robustness error family, conditional on supplied
  extraction data. Let $C\geq1$, $a>1$, and $0<b<1$. There are constants
  $a'\geq1$ and $0<b'<1$, depending only on $C,a,b$, such that for every
  admissible $(q,m,d)$, every $0\leq\eps\leq1$, every projective Pauli-test
  setting $S$, every global polynomial-pair witness $w$ for $S$ at the global
  error scale of Lemma~\ref{lem:qld-unitary}, and every extraction witness $v$
  for $S,w$ at the extraction error scale of that lemma, the state distance of
  the data transferred as in
  Theorem~\ref{thm:pauli-extraction-isometry-construction-support} is at most
  $a'(md)^{a'}(\eps^{b'}+q^{-b'}+2^{-b'md})$.
\end{example}
\begin{proof}
  \leanok
  \uses{lem:qld-extraction-error-form,
    lem:pauli-extraction-state-distance-support,
    thm:pauli-error-square-root-support}
  Take the constants supplied by the preservation of the robustness error
  family and halve the exponent. The squared state component gives the state
  distance as a square root, and the square-root estimate for the error family
  absorbs it into the displayed form.
\end{proof}

\begin{example}[Error-form isometry bounds for supplied extraction data]
  \label{lem:pauli-supplied-extraction-error-form-support}
  \lean{MIPStarRE.QPBT.pauli_soundness_deltaQld_ofExtractionWitness}
  \leanok
  \uses{lem:qld-unitary, lem:qld-extraction-error-form, def:povm-distance}
  This formalization-only node is conditional on supplied extraction data; it
  is not the soundness theorem of the paper. Let $C\geq1$, $a>1$, and
  $0<b<1$. There are constants $a'\geq1$ and $0<b'<1$, depending only on
  $C,a,b$, such that for every admissible $(q,m,d)$, every $0\leq\eps\leq1$,
  every projective Pauli-test setting $S$, every global polynomial-pair
  witness $w$ for $S$ at the global error scale of
  Lemma~\ref{lem:qld-unitary}, and every extraction witness $v$ for $S,w$ at
  the extraction error scale of that lemma, there are local isometries and a
  normalized auxiliary state whose state norm distance and whose two complete
  squared Pauli-family distances, for every $W\in\{X,Z\}$ and summed over all
  $h\in\F_q^M$, are all at most
  $a'(md)^{a'}(\eps^{b'}+q^{-b'}+2^{-b'md})$. Neither the global
  polynomial-pair measurement nor the extraction data is constructed here.
\end{example}
\begin{proof}
  \leanok
  \uses{thm:pauli-concrete-isometry-transfer-support,
    lem:qld-extraction-error-form,
    thm:pauli-error-square-root-support,
    lem:delta-qld-mono-support}
  Take the constants supplied by the preservation of the robustness error
  family, multiply the prefactor by four, and halve the exponent. Apply the
  concrete transfer to the supplied data. The square-root estimate for the
  error family bounds the state term, monotonicity in the universal constants
  bounds the two operator terms after multiplication by four, and the enlarged
  prefactor absorbs that factor.
\end{proof}

\begin{theorem}[Pauli isometry bounds for projective settings]
  \label{thm:pauli-projective-setting-isometry-support}
  \lean{MIPStarRE.QPBT.exists_projective_setting_isometry_bounds}
  \leanok
  \uses{def:strategy-observables, def:povm-distance, def:expanded-state}
  This formalization-only result concerns the projective setting in the
  appendix proof of Theorem~\ref{thm:pauli}, not the arbitrary strategies in
  its statement. There are universal constants $a\geq1$ and $0<b<1$, chosen
  before the parameters, error, and setting, such that for every admissible
  $(q,m,d)$, every $0\leq\eps\leq1$, and every projective Pauli-test setting
  $S$ winning with probability at least $1-\eps$, there are local isometries
  and a unit auxiliary state
  with state norm distance and both complete squared Pauli-family distances,
  for every $W\in\{X,Z\}$ and summed over all $h\in\F_q^M$, bounded by
  $a(md)^a(\eps^b+q^{-b}+2^{-bmd})$. No global polynomial-pair measurement
  or extraction data are assumed. The extension to arbitrary successful
  strategies and the seed-indexed game correspondence remain open; see
  \cite{gap:qpbt_extraction-transfer}.
\end{theorem}
\begin{proof}
  \leanok
  \uses{lem:qld-unitary, lem:pauli-supplied-extraction-error-form-support}
  Fix the universal constants and the extraction constant of
  Lemma~\ref{lem:qld-unitary}, and with them the constants of the error-form
  bounds for supplied extraction data. For the given setting, that lemma
  produces a global polynomial-pair witness and an extraction witness at the
  corresponding error scales. Feeding them to the error-form bounds gives the
  isometries, the auxiliary state, and the common bound in the displayed
  robustness family.
\end{proof}

\begin{theorem}[Questionwise Pauli Naimark reduction]
  \label{thm:pauli-naimark-reduction-support}
  \lean{MIPStarRE.QPBT.pauliNaimarkEmbeddingA,
    MIPStarRE.QPBT.pauliNaimarkEmbeddingB,
    MIPStarRE.QPBT.pauliNaimarkStrategy,
    MIPStarRE.QPBT.pauliNaimarkStrategy_isProjective,
    MIPStarRE.QPBT.pauliNaimarkStrategy_value,
    MIPStarRE.QPBT.pauliNaimarkStrategy_state,
    MIPStarRE.QPBT.pauliNaimarkSetting,
    MIPStarRE.QPBT.pauliNaimarkStrategy_compressA,
    MIPStarRE.QPBT.pauliNaimarkStrategy_compressB,
    MIPStarRE.QPBT.pauliNaimarkStrategy_pauli_compressA,
    MIPStarRE.QPBT.pauliNaimarkStrategy_pauli_compressB,
    MIPStarRE.QPBT.pauliNaimarkWitness,
    MIPStarRE.QPBT.pauli_naimark_witness_state_eq,
    MIPStarRE.QPBT.pauli_naimark_witness_state_distance_eq}
  \leanok
  \uses{thm:naimark}
  Let $S$ be any finite-dimensional strategy for the Pauli basis test,
  with possibly distinct local spaces. There is a projective strategy $T$
  on separately enlarged local spaces with $\operatorname{val}(T)
  =\operatorname{val}(S)$, and fixed question-independent local isometries
  $i_{\mathrm A},i_{\mathrm B}$ such that
  $\ket{\psi_T}=(i_{\mathrm A}\otimes i_{\mathrm B})\ket{\psi_S}$.
  Compression to the ground coordinates recovers both entire original
  question-indexed measurement families, including their full
  $\mathsf{Pauli}$-answer postprocessings (which fold wrong-form answers into
  the zero outcome). Thus $\operatorname{val}(S)\geq1-\eps$ constructs a
  projective setting for $T$ at the same error. For any soundness witness on
  $T$, composing its two local isometries with $i_{\mathrm A}$ and
  $i_{\mathrm B}$ gives a witness on $S$ with the same auxiliary state:
  the extracted state vectors, and hence their distances from the common
  ideal state, agree exactly. This is formalization-only support for the
  dilation at the start of the proof of Theorem~\ref{thm:pauli}; it does not
  identify compression with operator conjugation, transfer either complete
  operator-family estimate, or handle the unrestricted error domain.
\end{theorem}
\begin{proof}
  \leanok
  \uses{thm:naimark}
  Apply the completed one-measurement Naimark dilation for every question,
  with one fixed ground coordinate on each player's enlarged local space.
  Ground-slice compression recovers every POVM effect, and commutes with
  answer postprocessing. The padded state is the tensor image of the two
  ground embeddings, so all game correlations and the value are unchanged.
  Composition of local isometries gives equality of extracted states for
  the original and dilated strategies.
\end{proof}

\begin{theorem}[Pauli operator-family transfer through Naimark compression]
  \label{thm:pauli-naimark-operator-transfer-support}
  \lean{MIPStarRE.QPBT.pauli_naimark_operator_distanceA_le,
    MIPStarRE.QPBT.pauli_naimark_operator_distanceB_le}
  \leanok
  \uses{thm:pauli-naimark-reduction-support,
    thm:pauli-range-projection-support}
  Let $S$ be any strategy for the Pauli basis test, let $T$ be its
  questionwise Naimark dilation, and let $w$ be a soundness witness for $T$.
  For every $W\in\{X,Z\}$, write
  \[
    \Delta = (\phi_{\mathrm A}\ot\phi_{\mathrm B})\ket{\psi_T}
      \quad\text{and}\quad
    \Theta = \ket{\mathrm{aux}}\ot\ket{\mathrm{EPR}_q}^{\ot M}.
  \]
  The witness obtained by composing its local isometries with the two ground
  embeddings satisfies
  \begin{align*}
    d_{\mathrm A}(S) &\leq 3d_{\mathrm A}(T)+6\|\Delta-\Theta\|^2,\\
    d_{\mathrm B}(S) &\leq 3d_{\mathrm B}(T)+6\|\Delta-\Theta\|^2.
  \end{align*}
  Here each $d_w$ is the complete Pauli operator-family distance, summed over
  all answers.  This is a formalization-only support result: it has no
  outcome-cardinality factor and adds no consistency, projectivity, range, or
  producer hypothesis to the source theorem.
\end{theorem}
\begin{proof}
  \leanok
  \uses{thm:pauli-naimark-reduction-support,
    thm:pauli-range-projection-support}
  The range projection of each composed ground embedding fixes the dilated
  state.  The complete Pauli projector family on the opposite register gives
  the required mirror identity on the ideal EPR state.  The range-projection
  estimate therefore bounds the compressed family by three times the dilated
  family distance and six times the squared state distance.  The two final
  equalities are the exact compression identities for the pulled-back witness.
\end{proof}

\begin{theorem}[Pauli isometry bounds for arbitrary strategies]
  \label{thm:pauli-arbitrary-strategy-isometry-support}
  \lean{MIPStarRE.QPBT.exists_arbitrary_strategy_isometry_bounds}
  \leanok
  \uses{def:strategy-observables, def:povm-distance, def:expanded-state,
    thm:pauli-projective-setting-isometry-support,
    thm:pauli-naimark-reduction-support,
    thm:pauli-naimark-operator-transfer-support,
    lem:delta-qld-scalar-absorption-support}
  This formalization-only result removes the projectivity restriction of
  Theorem~\ref{thm:pauli-projective-setting-isometry-support} and keeps the
  error restriction $0\leq\eps\leq1$.  There are universal constants $a\geq1$
  and $0<b<1$, chosen before the parameters, the error and the strategy, such
  that for every admissible $(q,m,d)$, every $0\leq\eps\leq1$ and every
  strategy for $\mathfrak{G}^{\mathrm{Pauli}}_{(q,m,d)}$ succeeding with
  probability at least $1-\eps$, there are local isometries and a unit
  auxiliary state whose state norm distance and whose two complete squared
  Pauli-family distances, for every $W\in\{X,Z\}$ and summed over all
  $h\in\F_q^M$, are bounded by $a(md)^a(\eps^b+q^{-b}+2^{-bmd})$.  Neither
  projectivity nor extraction data is assumed.
\end{theorem}
\begin{proof}
  \leanok
  \uses{thm:pauli-projective-setting-isometry-support,
    thm:pauli-naimark-reduction-support,
    thm:pauli-naimark-operator-transfer-support,
    lem:delta-qld-scalar-absorption-support,
    lem:delta-qld-mono-support}
  Dilate the strategy questionwise by the Naimark reduction.  The dilation is
  projective and has exactly the same value, so it is a projective setting at
  the same error, and the projective-setting bounds apply to it.  Pull the
  resulting witness back along the two ground embeddings.  The state distance
  is unchanged by that pullback, so monotonicity in the universal constants
  brings it to the enlarged scale.  Each operator-family distance is at most
  three times the dilated family distance plus six times the squared state
  distance; both of those are bounded by the same error, and the scalar
  absorption returns $3\delta+6\delta^2$ to the robustness form with prefactor
  $21a^2$ and the same exponent.
\end{proof}

We can now prove the soundness theorem for the Pauli basis test, stated as Theorem~\ref{thm:pauli} in Chapter~\ref{chap:qpbt-test}.

\begin{proof}
	\proves{thm:pauli}
	\uses{lem:qld-unitary, lem:qld-extraction-error-form, lem:pauli-completeness,
		thm:naimark, lem:projective-strategy-setup, def:expanded-state,
		def:EPR, def:povm-distance, thm:pauli-ancilla-intertwining-support,
		thm:pauli-naimark-reduction-support,
		thm:pauli-naimark-operator-transfer-support,
		thm:pauli-isometry-transfer-alice-support,
		thm:pauli-isometry-transfer-bob-support}
	The account below is the route the source prints.  The Lean proof of
	Theorem~\ref{thm:pauli} performs the same Naimark reduction and is
	assembled from the formalization-support node
	Theorem~\ref{thm:pauli-arbitrary-strategy-isometry-support} together with
	the extension of the error domain past one; it is complete under the
	standard axioms, which is what the statement-level formalization mark on
	Theorem~\ref{thm:pauli} records.  No proof-level mark is placed on this
	account, because the divisibility obligation disclosed below is neither
	discharged nor altered by that closure.
	Let $(\ket{\psi},M)$ be a strategy for $\mathfrak{G}^{\mathrm{Pauli}}_{(q,m,d)}$ that succeeds with probability at least $1 - \eps$, where $\ket{\psi}$ is a bipartite state on registers $\mathrm{A}$ and $\mathrm{B}$. By Naimark's theorem (Theorem~\ref{thm:naimark}, as formulated in~\cite[Theorem 5.1]{Ji2020LowIndividualDegree}), at the cost of extending $\ket{\psi}$ by sufficiently many ancilla qubits initialized in the state $\ket{0}$, we may assume that the strategy is projective and that the state is padded as in the conventions of Lemma~\ref{lem:projective-strategy-setup} and the convention following it. Let $V_{\mathrm{A}}$ be the unitary acting on registers $\mathrm{A}\mathrm{A}'\mathrm{A}''$ given by Lemma~\ref{lem:qld-unitary}, and let $\phi_{\mathrm{A}}$ be the isometry mapping register $\mathrm{A}$ to registers $\mathrm{A}\mathrm{A}'\mathrm{A}''$, where $\mathrm{A}'$ and $\mathrm{A}''$ are isomorphic to $(\C^q)^{\ot M}$, defined by
	\[
		\phi_{\mathrm{A}}\colon \ket{\theta}_{\mathrm{A}} \mapsto V_{\mathrm{A}} \bigl(\ket{\theta}_{\mathrm{A}} \ot (\ket{\mathrm{EPR}_q}^{\ot M})_{\mathrm{A}' \mathrm{A}''}\bigr)\;;
	\]
	define $\phi_{\mathrm{B}}$ analogously.
	Theorem~\ref{thm:pauli-ancilla-intertwining-support} gives the isometry and
	exact intertwining identities for these unit-ancilla embeddings.
	% Reliance disclosure added: the proof chain of thm:pauli passes through
	% lem:qld-4-7, whose source proof (lines 1277--1413 of
	% 14_analysis_of_the_pauli_basis_test.tex) instantiates the classical low
	% individual degree test at dimension 2m+2 without a divisibility
	% hypothesis; the sentence below records the resulting outstanding
	% obligation.
	This invocation is where the analysis of the classical low individual degree
	test enters. Lemma~\ref{lem:qld-unitary} rests on the polynomial-pair
	measurement of Lemma~\ref{lem:qld-4-7}, whose proof instantiates the
	classical test in $2m+2$ variables and at present does so only when $2m+2$
	divides $q$ --- a divisibility that fails for every admissible tuple with
	$m \geq 2$. Three further steps of that same chain are likewise unproved.
	The classical soundness statement invoked there,
	Theorem~\ref{lem:ld-soundness},
	is imported from Theorem~4.7 of~\cite{Ji2021Quantum} through an
	identification of the game $\mathfrak{G}^{\mathrm{ld}}_{(q,m,d,1)}$ with
	the two-prover tensor-code test for $\code^{\ot m}$, and neither that game
	correspondence nor the bound on the auxiliary parameter of Theorem~4.7 under
	which the stated error function $\delta_{\mathrm{ld}}$ is obtained has been
	established. Both, with the divisibility obstruction, are documented in
	\cite{gap:qpbt_ld-dimension-divisibility}. Furthermore, the combined-lines
	consistency of Lemma~\ref{lem:qld-4-13}, on which
	Lemma~\ref{lem:qld-4-7} relies, is established here only with the error
	$m \cdot \poly(\eps, md/q)$, the error form $\poly(m^2 \eps, md/q)$
	carried in its statement remaining unproved; the discrepancy and the
	established form are documented in
	\cite{gap:qpbt_combined-lines-error-term}. The present proof is therefore
	complete only modulo these four obligations. Since the
	expanded state $\ket{\hat\psi}$ of Definition~\ref{def:expanded-state} is
	exactly $\ket{\psi}$ with the halves of $M$ pairs
	$\ket{\mathrm{EPR}_q}$ appended locally in registers
	$\mathrm{A}'\mathrm{A}''$ and $\mathrm{B}'\mathrm{B}''$, we have
	$\phi_{\mathrm{A}} \ot \phi_{\mathrm{B}} \ket{\psi}
	= V_{\mathrm{A}} \ot V_{\mathrm{B}} \ket{\hat\psi}$.
	Lemma~\ref{lem:qld-unitary} therefore yields a state $\ket{\mathrm{aux}}$
	on registers $\mathrm{A}\mathrm{A}'\mathrm{B}\mathrm{B}'$ such that
	\[
		\bigl\| \phi_{\mathrm{A}} \ot \phi_{\mathrm{B}} \ket{\psi} - \ket{\mathrm{aux}} \ot \ket{\mathrm{EPR}_q}^{\ot M} \bigr\|^2 \leq \delta_{\mathrm{qld}}\;.
	\]
	% Transfer step added: the source (lines 1864--1875 of
	% 14_analysis_of_the_pauli_basis_test.tex) passes directly from the
	% conjugation estimates of lem:qld-unitary, which concern the
	% unitarily conjugated operators V (M_h ot Id) V^dagger, to the
	% corresponding estimates for the isometry phi, whose conjugation
	% inserts the projection onto its range; the range-projection
	% argument below supplies the missing step.
	Item~2 of Lemma~\ref{lem:qld-unitary} concerns the unitarily conjugated operators $N^W_h = V_{\mathrm{A}} \, (M^{(\mathrm{Pauli},W)}_h)_{\mathrm{A}} \, V_{\mathrm{A}}^\dagger$, whereas conjugation by the isometry $\phi_{\mathrm{A}}$ additionally inserts the projection $Q_{\mathrm{A}} = \phi_{\mathrm{A}} \phi_{\mathrm{A}}^\dagger = V_{\mathrm{A}} \, \bigl(\ket{\mathrm{EPR}_q}\bra{\mathrm{EPR}_q}^{\ot M}\bigr)_{\mathrm{A}'\mathrm{A}''} \, V_{\mathrm{A}}^\dagger$ onto the range of $\phi_{\mathrm{A}}$: since $(M^{(\mathrm{Pauli},W)}_h)_{\mathrm{A}}$ and $\bigl(\ket{\mathrm{EPR}_q}\bra{\mathrm{EPR}_q}^{\ot M}\bigr)_{\mathrm{A}'\mathrm{A}''}$ act on disjoint registers and hence commute,
	\[
		\phi_{\mathrm{A}} \, (M^{(\mathrm{Pauli},W)}_h)_{\mathrm{A}} \, \phi_{\mathrm{A}}^\dagger
		= V_{\mathrm{A}} \, \Bigl( (M^{(\mathrm{Pauli},W)}_h)_{\mathrm{A}} \ot \bigl(\ket{\mathrm{EPR}_q}\bra{\mathrm{EPR}_q}^{\ot M}\bigr)_{\mathrm{A}'\mathrm{A}''} \Bigr) \, V_{\mathrm{A}}^\dagger
		= N^W_h \, Q_{\mathrm{A}} = Q_{\mathrm{A}} \, N^W_h\;.
	\]
	The estimates transfer from $N^W_h$ to $\phi_{\mathrm{A}} \, (M^{(\mathrm{Pauli},W)}_h)_{\mathrm{A}} \, \phi_{\mathrm{A}}^\dagger$ by Theorem~\ref{thm:pauli-isometry-transfer-alice-support}, whose range-projection step is as follows. Write $\ket{\Delta} = \ket{\mathrm{aux}} \ot \ket{\mathrm{EPR}_q}^{\ot M}$ and $\ket{\vartheta} = \phi_{\mathrm{A}} \ot \phi_{\mathrm{B}} \ket{\psi}$. Since $\phi_{\mathrm{A}}^\dagger \phi_{\mathrm{A}} = \Id$, the projection $\Id - Q_{\mathrm{A}}$ annihilates $\ket{\vartheta}$, so the displayed bound above gives
	\[
		\bigl\| (\Id - Q_{\mathrm{A}}) \ket{\Delta} \bigr\|^2
		= \bigl\| (\Id - Q_{\mathrm{A}}) (\ket{\Delta} - \ket{\vartheta}) \bigr\|^2
		\leq \bigl\| \ket{\Delta} - \ket{\vartheta} \bigr\|^2
		\leq \delta_{\mathrm{qld}}\;.
	\]
	Since the strategy is projective, the operators $N^W_h$, for $h \in \F_q^M$, are mutually orthogonal projectors summing to the identity, so $\sum_h \| N^W_h \ket{\xi} \|^2 = \| \ket{\xi} \|^2$ for every vector $\ket{\xi}$. Decomposing
	\[
		\bigl( \phi_{\mathrm{A}} \, (M^{(\mathrm{Pauli},W)}_h)_{\mathrm{A}} \, \phi_{\mathrm{A}}^\dagger - (\tau^W_h)_{\mathrm{A}''} \bigr) \ket{\Delta}
		= \bigl( N^W_h - (\tau^W_h)_{\mathrm{A}''} \bigr) \ket{\Delta} - N^W_h (\Id - Q_{\mathrm{A}}) \ket{\Delta}
	\]
	and summing the squared norms over $h$ via $\| \xi - \xi' \|^2 \leq 2\|\xi\|^2 + 2\|\xi'\|^2$, the first terms contribute $O(\delta_{\mathrm{qld}})$ by Item~2 of Lemma~\ref{lem:qld-unitary} and the second terms contribute at most $2 \| (\Id - Q_{\mathrm{A}}) \ket{\Delta} \|^2 \leq 2\delta_{\mathrm{qld}}$; hence, for $W \in \{X,Z\}$,
	\[
		\phi_{\mathrm{A}} \, (M^{(\mathrm{Pauli},W)}_h)_{\mathrm{A}} \, \phi_{\mathrm{A}}^\dagger \approx_{\delta_{\mathrm{qld}}} (\tau^W_h)_{\mathrm{A}''}
		\qquad \text{and, by the identical argument,} \qquad
		\phi_{\mathrm{B}} \, (M^{(\mathrm{Pauli},W)}_h)_{\mathrm{B}} \, \phi_{\mathrm{B}}^\dagger \approx_{\delta_{\mathrm{qld}}} (\tau^W_h)_{\mathrm{B}''}\;,
	\]
	where the $\approx$ statements hold with respect to the state $\ket{\mathrm{aux}} \ot \ket{\mathrm{EPR}_q}^{\ot M}$ and the answer summation is over $h \in \F_q^M$; the constant implicit in $\approx_{\delta_{\mathrm{qld}}}$ (Definition~\ref{def:povm-distance}) absorbs the factors accumulated above. The Bob comparison is the corresponding application of Theorem~\ref{thm:pauli-isometry-transfer-bob-support}. This shows that the game $\mathfrak{G}^{\mathrm{Pauli}}_{(q,m,d)}$ is a self-test, with robustness $\delta_{\mathrm{qld}}(\eps,m,d,q)$, for the canonical Pauli strategy of Lemma~\ref{lem:pauli-completeness}. Identifying the registers $\mathrm{A}\mathrm{A}'$ with $\mathrm{A}'$ and $\mathrm{B}\mathrm{B}'$ with $\mathrm{B}'$ matches the statement of Theorem~\ref{thm:pauli}, in which $\phi_{\mathrm{A}}\colon \mathcal{H}_{\mathrm{A}} \to \mathcal{H}_{\mathrm{A}'} \ot \mathcal{H}_{\mathrm{A}''}$. Finally, unpacking the dependence on the parameters $\eps, m, d, q$, the function $\delta_{\mathrm{qld}}(\eps,m,d,q) = O(\delta_P^{1/4} + md/q)$ has the form $a(md)^a(\eps^b + q^{-b} + 2^{-bmd})$ for universal constants $a > 1$ and $0 < b < 1$; here $\delta_P = C(\delta_G + \sqrt{\eps} + md/q)$ and $\delta_G$ already has that form. See Remark~\ref{rem:pauli-robustness-form} for the reconciliation with the statement of Theorem~\ref{thm:pauli}.
	The range-projection transfer just used, which is omitted in the source,
	is documented in \cite{gap:qpbt_extraction-transfer}.
\end{proof}

\begin{remark}[Norm conventions in the robustness bound]\label{rem:pauli-robustness-form}
	The proof above, following the source, establishes the first conclusion of Theorem~\ref{thm:pauli} in the squared form $\| \phi_{\mathrm{A}} \ot \phi_{\mathrm{B}} \ket{\psi} - \ket{\mathrm{aux}} \ot \ket{\mathrm{EPR}_q}^{\ot M} \|^2 \leq \delta_{\mathrm{qld}}$, whereas the theorem bounds the unsquared norm by $\delta_{\mathrm{qld}}(\eps,m,d,q)$. The two are reconciled by replacing the exponent $b$ with $b/2$ and enlarging $a$, since the class of functions $a(md)^a(\eps^b + q^{-b} + 2^{-bmd})$ is closed under square roots up to such adjustments of the universal constants. Similarly, the closing sentence of the source proof asserts constants $a > 1$ and $0 < b < 1$, while Theorem~\ref{thm:pauli} requires only $a \geq 1$ and $0 < b < 1$; the former implies the latter.
\end{remark}

\section{Formalization support for the robustness bound}

The following monotonicity statement is not a named result of the source
article.  It records, as a single inequality, the adjustment of the universal
constants of Theorem~\ref{thm:pauli} that is used when the robustness bound is
passed from its squared to its unsquared form, and when several estimates are
brought to a common scale.

\begin{lemma}[Monotonicity of the robustness bound in its universal constants]\label{lem:delta-qld-mono-support}
	\lean{MIPStarRE.QPBT.deltaQld_mono}
	\leanok
	\uses{thm:pauli}
	Write $\delta_{\mathrm{qld}}^{a,b}(\eps,m,d,q) = a(md)^a\bigl(\eps^b + q^{-b} + 2^{-bmd}\bigr)$ for the soundness function of Theorem~\ref{thm:pauli} with universal constants $a$ and $b$. Let $(q,m,d)$ be an admissible parameter tuple (Definition~\ref{def:admissible}), let $1 \le a \le a'$, let $0 < b' \le b$, and let $0 \le \eps \le 1$. Then
	\[
		\delta_{\mathrm{qld}}^{a,b}(\eps,m,d,q) \le \delta_{\mathrm{qld}}^{a',b'}(\eps,m,d,q)\;.
	\]
\end{lemma}
\begin{proof}
	\leanok
	\uses{def:admissible}
	Admissibility gives $md \ge 1$ and $q \ge 1$, so $a(md)^a \le a'(md)^{a'}$. In the bracket, $\eps^b \le \eps^{b'}$ because $0 \le \eps \le 1$ and $0 < b' \le b$; $q^{-b} \le q^{-b'}$ because $q \ge 1$ and $-b \le -b'$; and $2^{-bmd} \le 2^{-b'md}$ because $md \ge 0$ and $-bmd \le -b'md$. Both factors are nonnegative, so the two products compare factorwise.
\end{proof}

The next statement is likewise not a named result of the source article.  It
records the constant bookkeeping that returns the operator-transfer bound of
the Naimark reduction, which has the shape $3\delta + 6\delta^2$, to the
robustness form of Theorem~\ref{thm:pauli}.  The source performs the same
absorption inline when it records that $\delta_{\mathrm{qld}}$ has the stated
form.

\begin{lemma}[Scalar absorption of the transfer constants]\label{lem:delta-qld-scalar-absorption-support}
	\lean{MIPStarRE.QPBT.three_add_six_sq_deltaQld_le}
	\leanok
	\uses{thm:pauli, lem:delta-qld-mono-support}
	Write $\delta_{\mathrm{qld}}^{a,b}(\eps,m,d,q) = a(md)^a\bigl(\eps^b + q^{-b} + 2^{-bmd}\bigr)$ as in Lemma~\ref{lem:delta-qld-mono-support}, and abbreviate $\delta = \delta_{\mathrm{qld}}^{a,b}(\eps,m,d,q)$.  Let $(q,m,d)$ be an admissible parameter tuple (Definition~\ref{def:admissible}), let $a \ge 1$, let $b > 0$, and let $0 \le \eps \le 1$.  Then
	\[
		3\delta + 6\delta^2 \le \delta_{\mathrm{qld}}^{21a^2,\,b}(\eps,m,d,q)\;.
	\]
\end{lemma}
\begin{proof}
	\leanok
	\uses{def:admissible}
	Admissibility gives $md \ge 1$ and $q \ge 1$.  Since $0 \le \eps \le 1$ and $b > 0$ we have $\eps^b \le 1$; since $q \ge 1$ and $-b \le 0$ we have $q^{-b} \le 1$; and since $bmd \ge 0$ we have $2^{-bmd} \le 1$.  The bracket $T = \eps^b + q^{-b} + 2^{-bmd}$ therefore satisfies $0 \le T \le 3$.  Writing $K = a(md)^a$, admissibility and $a \ge 1$ give $K \ge 1$, and $\delta = KT$ satisfies $0 \le \delta \le 3K$.  Hence
	\[
		21K\delta - 3\delta - 6\delta^2 = 6\,\delta(3K - \delta) + 3\,\delta(K - 1) \ge 0,
	\]
	so $3\delta + 6\delta^2 \le 21K\delta = 21a^2 (md)^{2a} T$.  Finally $md \ge 1$ and $2a \le 21a^2$ give $(md)^{2a} \le (md)^{21a^2}$, and the right-hand side is exactly $\delta_{\mathrm{qld}}^{21a^2,b}(\eps,m,d,q)$.
\end{proof}
